\documentclass{lmcs}

\usepackage{amsmath,amssymb}
\usepackage{hyperref}

\newcommand{\I}{\mathbb{I}}
\newcommand{\DM}[1]{\mathrm{DM}(#1)}
\newcommand{\transp}{\mathsf{transp}}
\newcommand{\hcomp}{\mathsf{hcomp}}
\newcommand{\vcomp}{\mathsf{comp}}
\newcommand{\refl}{\mathsf{refl}}
\newcommand{\J}{\mathsf{J}}
\newcommand{\ua}{\mathsf{ua}}
\newcommand{\Glue}{\mathsf{Glue}}
\newcommand{\glue}{\mathsf{glue}}
\newcommand{\unglue}{\mathsf{unglue}}
\newcommand{\nf}{\mathsf{nf}}
\newcommand{\FV}{\mathsf{FV}}
\renewcommand{\conv}{\equiv}
\newcommand{\aconv}{\equiv_{\mathrm{alg}}}
\newcommand{\repeq}{\approx}
\newcommand{\valeq}{\equiv^{\mathrm{val}}}
\newcommand{\Path}{\mathsf{Path}}
\providecommand{\fst}{}\renewcommand{\fst}{\mathsf{fst}}
\providecommand{\snd}{}\renewcommand{\snd}{\mathsf{snd}}
\newcommand{\U}{\mathcal{U}}
\newcommand{\ii}{\mathsf{i1}}
\newcommand{\io}{\mathsf{i0}}
\newcommand{\evalop}{\mathsf{eval}}
\newcommand{\constq}{\mathsf{const?}}
\newcommand{\constspec}{\mathsf{const?}_{\mathrm{spec}}}
\newcommand{\convop}{\mathsf{conv}}
\newcommand{\quo}{\mathsf{quote}}
\newcommand{\capp}{\mathsf{inst}}
\newcommand{\Comp}{\mathsf{Comp}}

\begin{document}

\title[Evaluation-Time Constancy Inference]{Evaluation-Time Constancy
  Inference for Transport in Cubical Type Theory}

\author[R.~Shibusawa]{Ryota Shibusawa}
\address{Daiichi Institute of Technology, Japan}
\email{r-shibusawa@daiichi-koudai.ac.jp}

\keywords{cubical type theory, definitional equality, regularity,
  transport, normalization by evaluation, De Morgan algebra}

% REMOVE BEFORE SUBMISSION: draft version marker (added so that
% successive external review rounds can identify which revision they
% are reading).
\titlecomment{Draft v12, 2026-07-17.}

\begin{abstract}
  The eliminator $\J$ of cubical type theories does not compute
  definitionally on $\refl$: the transport primitive's constancy
  formula is fixed when a term is written, so constancy arising later
  by substitution is never used.  Re-detecting it at evaluation time
  is an old informal idea; this paper gives it a metatheory.  First,
  a no-go theorem: combining the evaluation-time rule equationally
  with the structural transport rules derives a constant-system
  $\hcomp$-regularity principle --- closely related to regularity
  principles known to fail in standard cubical-sets models ---
  already at path types, before $\Glue$; and restricting the
  structural rules to non-constant families is not
  substitution-stable.  Second, the positive system: judgmental
  equality is defined as the algorithmic equality of a prioritized
  normalization-by-evaluation strategy whose constancy check is
  specified representation-independently (the fresh dimension must
  not occur in the read-back of the family's value) and implemented
  exactly; we prove the check semantically sound and the typed
  algorithmic equality an admissible conversion relation ---
  equivalence, congruence, weakening, substitution, context
  conversion, type preservation --- with termination and canonicity
  relative to an explicitly delimited base component, which we
  discharge for a $\Glue$-free, universe-free core fragment: there,
  including for $\J\,d\,\refl \conv d$ itself, the results are
  unconditional.  All concrete positive and negative
  object-language conversion witnesses, including the switchover
  experiments behind the no-go theorem, are machine-checked in a
  self-contained Lean~4 kernel; the metatheory itself is
  pen-and-paper.
\end{abstract}

\maketitle

\section{Introduction}
\label{sec:intro}

Cubical type theories \cite{CCHM,ABCFHL} compute: univalence and
higher inductive types have computational content, and canonicity and
normalization hold \cite{HuberNorm,SterlingAngiuli}.  The price is
paid at the eliminator $\J$ of Martin-L\"of's identity type, which in
cubical systems is a defined operation --- transport along a
connection square --- and no longer satisfies the definitional
computation rule $\J\,d\,\refl \conv d$.  In Cubical Agda
\cite{VezzosiMortbergAbel} the equation holds only propositionally
(\textsf{transportRefl}); library code is peppered with explicit
appeals to it, and the gap is a standing source of friction when
porting Book-HoTT developments \cite{UIPCubical}.

The mechanism is easy to isolate.  The transport primitive
$\transp^i\,A\,\varphi\,u$ of ABCFHL-style theories carries a
\emph{constancy formula} $\varphi$: a cofibration under which the
family $A$ is required to be judgmentally constant in $i$, with the
primitive rule $\transp^i A\,\ii\,u \conv u$.  The formula is fixed
when the term is \emph{written}.  Inside the definition of $\J$ the
transport is formed with $\varphi = \io$; substituting $p := \refl$
later makes the family judgmentally constant, but nothing revisits
$\varphi$.  The information is available at evaluation time and is
simply not used.

The obvious repair has circulated informally for years: let the
\emph{evaluator} re-detect constancy --- normalize the family, check
that the transport dimension does not occur, and if so return the
input.  Arend's \textsf{coe} computes on constant families over its
(non-cubical) interval \cite{Arend}.  For full cubical theories the
folklore assessment --- which has, to our knowledge, circulated only
informally --- locates the obstruction in the interaction with the
structural transport rules, above all at $\Glue$, where regularity
fails in the standard models \cite{SwanRegularity,SattlerNote}; no
published metatheoretic analysis of the idea for a univalent cubical
calculus appears to exist.
This paper provides that analysis: a no-go theorem locating the
obstruction precisely, and the metatheory of an operational
formulation that avoids it.

\subparagraph*{The no-go half.}  Write (R) for the evaluation-time
rule: $\transp^i A\,\varphi\,u \conv u$ whenever $A$ is judgmentally
$i$-constant, regardless of $\varphi$.  When the family is constant,
(R) and the \emph{structural} rule for the family's head former both
apply, and an equational theory containing both closes over the pair.
We machine-check the resulting \emph{switchover} equations former by
former (Section~\ref{sec:switchover}).  For $\Pi$, $\Sigma$ and the
argument-wise higher inductive constructors the structural contractum
collapses back to the input --- using, and in fact \emph{forcing},
the $\eta$~laws.  For path types it does not: the structural rule
\[
  \transp^i (\Path\,A\,a\,b)\,\varphi\,p \;\conv\;
  \langle j \rangle\, \vcomp^i\,A\,[\,j{=}0 \mapsto a,\ j{=}1 \mapsto
  b\,]\,(p\,j)
\]
leaves, at a constant family, an $\hcomp$ whose tubes are constant in
the fill dimension.  Equating that residue with $p$ is a
constant-system homogeneous-composition regularity principle --- an
equation of the kind excluded in the standard cubical-sets models
\cite{SwanRegularity,SattlerNote}, and one our kernel refutes
($\mathsf{trans}\,\refl\,\refl \not\aconv \refl$).  Thus
(Theorem~\ref{thm:nogo}, with the full derivation in
Section~\ref{sec:nogo}): \emph{any equational theory containing the
structural path-transport rule, rule (R), congruence and transitivity
derives constant-tube $\hcomp$-regularity; the conflict arises at
path types already --- before, and independently of, $\Glue$.}
Moreover the alternative of \emph{restricting} the structural rules
to non-constant families --- the route taken by an earlier version of
this work --- fails for a dual reason: non-constancy is \emph{not}
stable under substitution (a family non-constant at definition time
is exactly one that may become constant later), so the restricted
rule set is not closed under substitution either.

\subparagraph*{The positive half.}  The way out is to stop generating
the equality by rules on terms and instead take the judgmental
equality of the extended system to \emph{be} the algorithmic equality
of a prioritized evaluation strategy: normalization by evaluation, in
which a transport node first evaluates its family, instantiates it at
a fresh dimension, and checks constancy \emph{of the resulting
value}; only if the check fails does it dispatch structurally.  For
this to be a definition of judgmental equality rather than a
description of an implementation, a body of metatheory is owed, and
this paper's main technical contribution is to supply it
(Sections~\ref{sec:domain}--\ref{sec:canonicity}):

\begin{itemize}
\item a \emph{representation equivalence} $\repeq$ on the semantic
  domain (values differing only in closure representation), because
  in a defunctionalized NbE the two sides of the substitution lemma
  do \emph{not} produce identical values, only $\repeq$-related ones
  (Section~\ref{sec:domain});
\item a \emph{specification} of the constancy check ---
  the fresh dimension must not occur in the \emph{read-back} of the
  family value --- which is representation-independent, together with
  its \emph{semantic soundness theorem}: a checked family has
  logically-related endpoints (Section~\ref{sec:check});
  the implemented occurs-check is proved to \emph{decide the
  specification exactly} (Theorem~\ref{thm:exact}), so the kernel and
  the specified system are one system, not two;
\item the substitution lemma
  $\evalop(t\sigma, \rho) \repeq \evalop(t, \evalop(\sigma,\rho))$,
  whose transport case goes through \emph{because} the check is
  specified representation-independently --- the branch taken is a
  function of the $\repeq$-class (Section~\ref{sec:oper});
\item a Kripke logical relation $\valeq$ over the domain whose
  fundamental theorem yields, as corollaries, that the typed
  algorithmic equality is an \emph{admissible conversion relation}:
  an equivalence, a congruence for every term former, stable under
  weakening, substitution and context conversion, type-preserving,
  and coherent with the bidirectional checker
  (Theorem~\ref{thm:admissible});
\item termination folded into the logical relation in the standard
  Tait style --- no syntactic termination measure is claimed, and the
  double induction (typing derivation outside, semantic type values
  inside) discharging the structural transport branch is made
  explicit (Section~\ref{sec:canonicity});
\item canonicity: closed naturals evaluate to numerals, whence
  consistency and $\J\,d\,\refl \conv d$
  (Theorem~\ref{thm:canonicity}).  The base-rule cases of the
  fundamental theorem are inherited as an explicitly delimited
  \emph{base component} (Section~\ref{sec:assumptions});
\item an \emph{unconditional core}: for the minimal fragment
  $\Pi, \Sigma, \Path, \mathbb{N}$ with transport and homogeneous
  composition --- no $\Glue$, no universes, no HITs --- we discharge
  the base component ourselves, in full per-former detail
  (Appendix~\ref{app:core}), so that there admissibility, canonicity
  and the definitional $\J\,d\,\refl$ hold with no assumption at all
  (Theorem~\ref{thm:core}); the derived $\J$ and its computation on
  $\refl$ are spelled out as an explicit proposition
  (Proposition~\ref{prop:J}).  The extension of the core by $S^1$ is
  stated separately as a proposition with a proof sketch and
  artifact evidence, so the headline result is not entangled with
  higher-inductive obligations;
\end{itemize}

\subparagraph*{The three layers, once and for all.}  The paper's
positive results live at three clearly separated strengths, and we
keep the labels attached throughout:
\begin{itemize}
\item \emph{Unconditional theorem} (Theorem~\ref{thm:core},
  Appendix~\ref{app:core}): the $\Glue$-free, universe-free
  \emph{minimal core} $\Pi, \Sigma, \Path, \mathbb{N}$ with
  $\transp$ and $\hcomp$ --- large enough to define $\J$, so
  $\J\,d\,\refl \conv d$ is unconditional;
\item \emph{Relative theorem} (Theorems~\ref{thm:fundamental},
  \ref{thm:admissible}, \ref{thm:canonicity}): the full implemented
  De~Morgan calculus, relative to the delimited base component (B)
  of Section~\ref{sec:assumptions};
\item \emph{Artifact evidence}: $\Glue$, univalence, and the HIT
  roster, machine-checked in the kernel but not covered by the
  unconditional layer.
\end{itemize}
When we speak of a metatheory ``for a univalent cubical calculus'',
univalence lives in the second and third layers; the unconditional
layer is deliberately $\Glue$- and universe-free.

We claim neither full regularity (the catalogue records
$\mathsf{trans}\,p\,\refl \not\aconv p$: $\hcomp$ is untouched, as by
Theorem~\ref{thm:nogo} it must be), nor a semantic countermodel for
the no-go equation (we state precisely what is refuted
algorithmically and what is conjectured semantically), nor a model
construction for the positive system; open problems are collected in
Section~\ref{sec:conclusion}.  A complementary result --- on a
generic path, algorithmic equality of connection-reparametrizations
is exactly free De~Morgan equality --- is kept in
Appendix~\ref{app:reparam}.

\section{The base calculus}
\label{sec:base}

\subsection{Syntax and rules}
\label{sec:rules}

We fix an ABCFHL-style univalent cubical type theory $T_0$ with a
De~Morgan interval.\footnote{The Cartesian variant carries the
switchover analysis through unchanged --- the residues are about
$\hcomp$, not connections --- and strengthens the metatheory
(Theorem A below).  We fix the De~Morgan case as primary because the
accompanying implementation realizes it.}  The judgment forms,
formers and rules used in this paper:

\begin{itemize}
\item \emph{Interval and cofibrations.}  Interval expressions over
  $i_1,\dots,i_n$ form the free De~Morgan algebra
  $\DM{i_1,\dots,i_n}$; cofibrations $\varphi$ are generated by face
  equations $i{=}0$, $i{=}1$, conjunction and disjunction, with
  constants $\io, \ii$.
\item \emph{Formers.}  $\Pi$, $\Sigma$, $\Path$ (heterogeneous
  $\mathsf{PathP}$), universes \`a la Russell, natural numbers and
  base types, $\Glue$, and the higher inductive types listed in
  Section~\ref{sec:hits}.
\item \emph{$\eta$~rules} for $\Pi$, $\Sigma$ and $\Path$.
\item \emph{Transport}, with an explicit constancy formula:
  \[
    \frac{\Gamma, i : \I \vdash A\ \mathsf{type} \qquad
          \Gamma \vdash \varphi\ \mathsf{cof} \qquad
          \Gamma, \varphi, i : \I \vdash A \conv A(\io/i) \qquad
          \Gamma \vdash u : A(\io/i)}
         {\Gamma \vdash \transp^i A\,\varphi\,u : A(\ii/i)}
  \]
  with the primitive collapse rule
  $(\mathrm{S}_{\ii})\colon \transp^i A\,\ii\,u \conv u$ and one
  \emph{structural} rule per former; the three needed later
  ($\varphi$ omitted where it merely distributes):
  \begin{align*}
    \transp^i (\Pi x{:}A.B)\,u &\conv
      \lambda x_1.\, \transp^i B[w\,x_1\,i/x]\,
        \bigl(u\,(w\,x_1\,0)\bigr),
      && w\,x_1\,i := \transp^j A(i \vee \neg j)\,x_1;\\
    \transp^i (\Sigma x{:}A.B)\,u &\conv
      \bigl(\transp^i A\,(\fst u),\
            \transp^i B[\overline{u}\,i/x]\,(\snd u)\bigr),
      && \overline{u}\,i := \transp^j A(i \wedge j)\,(\fst u);\\
    \transp^i (\Path\,A\,a\,b)\,p &\conv
      \langle j \rangle\, \vcomp^i A\,
        [\,j{=}0 \mapsto a,\ j{=}1 \mapsto b\,]\,(p\,j).
  \end{align*}
  Higher inductive types transport argument-wise on constructor
  arguments (Section~\ref{sec:hits}); $\Glue$ transports by the
  composite operation $\mathsf{transpGlue}$ of \cite[\S6.2]{CCHM},
  recalled in Section~\ref{sec:glueresidue}.
\item \emph{Homogeneous composition} $\hcomp$ with the usual boundary
  rules.  \emph{No} regularity rule for $\hcomp$ is assumed, and none
  may be: constant-tube collapse is of the kind excluded in the
  standard models \cite{SwanRegularity,SattlerNote}.  $\vcomp$
  abbreviates transport followed by $\hcomp$ of the transported
  tubes.
\end{itemize}

\subsection{The base component, and the three-layer main theorem}
\label{sec:assumptions}

The metatheory below is organized so that everything specific to rule
(R) is proved here, while the (large) base-theory content enters as
one delimited component:

\begin{itemize}
\item[(B)] \emph{The base component}: the Kripke logical relation of
  Section~\ref{sec:domain} restricted to $T_0$, i.e.\ the predicate
  clauses for all $T_0$ formers, the read-back correctness lemmas
  ($v \valeq \evalop(\quo\,v)$, uniqueness and inconvertibility of
  distinct normal forms, injectivity of type constructors), and the
  fundamental-theorem cases for all $T_0$ rules, in environment form.
\end{itemize}

\noindent Our results then come in three layers:

\begin{itemize}
\item \textbf{Theorem A$'$ (unconditional minimal core).}  For the
  fragment $\Pi, \Sigma, \Path, \mathbb{N}$ with $\transp$ and
  $\hcomp$ --- no $\Glue$, no universes, no HITs --- component (B)
  is discharged \emph{in this paper} (Appendix~\ref{app:core}), in
  per-former detail: because the fragment has no universe, its
  logical relation is defined by \emph{structural recursion on
  syntactic types in environments} --- a plainly well-founded
  induction --- and because it has no $\Glue$, transport never meets
  the $\mathsf{transpGlue}$ machinery.  All positive results are
  unconditional there (Theorem~\ref{thm:core}); in particular
  $\J\,d\,\refl \conv d$, whose definition needs only paths,
  connections and transport, is unconditional
  (Proposition~\ref{prop:J}).  The $S^1$ extension is a separate
  proposition (Appendix~\ref{app:core}, last subsection).
\item \textbf{Theorem A (Cartesian layer).}  For Cartesian cubical
  type theory, the judgmental-theory-level part of component (B) ---
  the normal-form language and bijection, inconvertibility of
  distinct normal forms, injectivity of type constructors,
  decidability --- is discharged by the normalization theorem of
  Sterling--Angiuli \cite{SterlingAngiuli} (normal-form specification
  in \cite{HuangNF}), \emph{for the fragment their theorem covers}:
  $\Pi$, $\Sigma$, path, glue and the circle, \emph{without
  universes}.  The evaluator-level part of (B) --- the concrete
  Kripke relation over our domain and the environment-form
  fundamental cases --- is not literally in \cite{SterlingAngiuli},
  whose method is deliberately not evaluation-based; it is a standard
  NbE-correctness transposition, which we enumerate but do not claim
  from the literature.  Appendix~\ref{app:bridge} gives the precise
  item-by-item correspondence, the translation between their
  boundary-conditioned $\mathsf{coe}$ and our $\varphi$-annotated
  $\transp$, and the exact residual gaps (universes; the wider HIT
  roster; the evaluator-level transposition).  Theorem A is therefore
  unconditional only over that fragment and modulo that
  transposition; we do not claim more.
\item \textbf{Theorem B (De~Morgan layer).}  For the De~Morgan
  calculus --- which the accompanying implementation realizes ---
  component (B) is an assumption: Huber \cite{HuberNorm} provides its
  closed/canonicity fragment, but a full open-term normalization
  proof on the De~Morgan side is not in the literature.  All theorems
  below are stated relative to (B).
\end{itemize}

\section{Rule (R), and the instability of non-constancy}
\label{sec:rule}

\begin{defi}[Judgmental constancy]\label{def:const}
  $A$ is \emph{judgmentally $i$-constant} (over $\Gamma, i : \I$) if
  there is $A_0$ with $\Gamma, i \vdash A \conv A_0$ and
  $i \notin \FV(A_0)$.
\end{defi}

\noindent The rule under study is
\[
  (\mathrm{R})\qquad
  \transp^i A\,\varphi\,u \conv u
  \qquad\text{for \emph{any} $\varphi$, whenever $A$ is judgmentally
  $i$-constant.}
\]
Write $T_{\mathrm{eq}} := T_0 + (\mathrm{R})$ for the
\emph{unrestricted equational theory}.  Two different systems appear
in this paper, and we name them apart once and for all:
$T_{\mathrm{eq}}$ is the object of the no-go analysis
(Section~\ref{sec:switchover}) and the target of the soundness
theorem (Section~\ref{sec:soundness}); it is not used as the type
system of record (its judgmental equality appears in statements, its
typing never does).
The system of record --- in which every fundamental-theorem,
admissibility and canonicity statement below is made --- is the type
system $T_{\mathrm{alg}}$ of Definition~\ref{def:talg}, whose
conversion rule uses the algorithmic equality.

\begin{rem}[Why the explicit formula matters]\label{rem:explicit}
  In a $\varphi$-free presentation whose collapse rule is the side
  condition ``$i \notin \FV(A) \Rightarrow \transp^i A\,u \conv u$'',
  rule (R) is derivable from congruence.  With the explicit formula
  --- as in ABCFHL and Cubical Agda --- it is not:
  $(\mathrm{S}_{\ii})$ applies only at $\varphi = \ii$, and
  congruence cannot change $\varphi$.  So (R) genuinely extends the
  judgmental equality of the $\varphi$-explicit calculus; for the
  $\varphi$-free reading, the contribution of this paper is the
  operational system of Section~\ref{sec:oper}.
\end{rem}

\begin{lem}[Positive substitution stability]\label{lem:posstab}
  If $A$ is judgmentally $i$-constant over $\Gamma, i$ and
  $\sigma : \Delta \to \Gamma$, then $A\sigma$ is judgmentally
  $i$-constant over $\Delta, i$.
\end{lem}
\begin{proof}
  Judgmental equality is closed under substitution; $i$ is bound at
  the transport site, may be chosen fresh for $\Delta$, and the terms
  of $\sigma$ live over $\Delta$, so $\FV(A_0\sigma) \not\ni i$.
\end{proof}

\begin{lem}[Type preservation]\label{lem:preservation}
  If $A$ is judgmentally $i$-constant then
  $A(\io) \conv A_0 \conv A(\ii)$; rule (R) preserves typing.
\end{lem}

\noindent Note the asymmetry: constancy is substitution-stable;
\emph{non}-constancy is not --- substituting $\refl$ for a path
variable is exactly how a non-constant family becomes constant.  Any
side condition of the form ``the family is \emph{not} constant'' is
therefore unusable in a rule schema, both because it is not preserved
by substitution and because, as a negative premise, it does not even
define an inductive rule set.

\section{The switchover analysis and the no-go theorem}
\label{sec:switchover}

When the family is constant, (R) and the structural rule for its head
former both apply.  Call the induced equation between the two
contracta the \emph{switchover equation} of the former.  This section
determines, former by former and machine-checked throughout, which
switchover equations hold in the algorithmic equality
(Section~\ref{sec:oper}); the probes are the file
\texttt{LibSwitchover.lean} of the artifact.

\subsection{Convergent formers}
\label{sec:convergent}

\begin{lem}[Switchover convergence at $\Pi$, $\Sigma$]
  \label{lem:converge}
  Let the families involved be judgmentally $i$-constant.  Then, in
  $T_0 + (\mathrm{R})$ and in the algorithmic equality:
  \begin{align*}
    \lambda x.\,\transp^i B\,\bigl(u\,(\transp^i A\,x)\bigr)
      &\conv u  &&(\Pi)\\
    \bigl(\transp^i A\,(\fst u),\ \transp^i B'\,(\snd u)\bigr)
      &\conv u  &&(\Sigma, \text{ incl.\ dependent})
  \end{align*}
\end{lem}
\begin{proof}
  $\Pi$: the inner families are reparametrizations of constant
  families, hence constant; (R) collapses both inner transports,
  leaving $\lambda x.\,u\,x$, and $\eta_\Pi$ finishes.  $\Sigma$: in
  the dependent case the second family is $B[\overline{u}\,i/x]$ with
  $\overline{u}\,i \conv \fst u$ by (R), hence constant; (R) twice,
  then $\eta_\Sigma$.  Machine witnesses: \texttt{swPiD},
  \texttt{swSigD}, \texttt{swSigDepD} (all accepted).
\end{proof}

\begin{rem}[(R) forces $\eta$]\label{rem:eta}
  Over an $\eta$-free base, the $\Sigma$ switchover equation \emph{is}
  surjective pairing: any theory containing (R) and structural
  $\Sigma$-transport derives $\eta_\Sigma$ (similarly $\eta_\Pi$).
  Our kernel has all three $\eta$~laws, machine-checked, as any
  implementation of (R) must.
\end{rem}

\subsection{Higher inductive types: the argument-wise class}
\label{sec:hits}

For higher inductive types we make no blanket claim; we delimit a
signature class and check it.

\begin{defi}[Argument-wise transport signatures]\label{def:argwise}
  A constructor signature is \emph{argument-wise} if its structural
  transport transports each point argument along the corresponding
  line and carries interval arguments unchanged, generating no
  boundary correction ($\hcomp$/filler).
\end{defi}

\begin{lem}\label{lem:hitconverge}
  For argument-wise signatures, the switchover equation converges by
  (R) on each argument line alone (no $\eta$ needed): the contractum
  is literally the constructor form.
\end{lem}

\noindent The signatures implemented in the artifact classify as
follows (machine witnesses in the last column):

\begin{center}
\small
\begin{tabular}{llll}
  HIT & constructors & class & switchover \\\hline
  lists & $\mathsf{nil}, \mathsf{cons}$ & argument-wise &
    converges (\texttt{swListD}) \\
  suspension & $\mathsf{N}, \mathsf{S}, \mathsf{merid}$ &
    arg.-wise (path ctor.) & converges (\texttt{swSuspMeridD}) \\
  pushout & $\mathsf{inl}, \mathsf{inr}, \mathsf{push}$ &
    argument-wise & converges (same schema) \\
  $S^1$, torus & point/loop/surface & no point arguments &
    trivially converges \\
  $K(G,1)$, $B\mathrm{Gpd}$ & base/loop/comp cells &
    argument-wise & converges (same schema) \\
  quotient & $\mathsf{qin}$ & argument-wise & converges \\
  quotient & $\mathsf{qeq}$, squash cells &
    \emph{no structural rule} & vacuous (no critical pair) \\
  truncation & $\mathsf{tin}$, squash & argument-wise / none &
    converges / vacuous \\
\end{tabular}
\end{center}

\noindent In our kernel, transports of non-constructor elements of
quotients and truncations become neutral, so for the squash-type
constructors there is no structural contractum and hence no
switchover question.  Signatures whose structural transport would
generate boundary corrections (constructors with non-trivial boundary
in the transported arguments) are \emph{not} covered by
Lemma~\ref{lem:hitconverge}; in the light of
Section~\ref{sec:nogo} we expect non-convergence there, and we leave
a general signature framework as future work.

\subsection{Non-convergence at path types}
\label{sec:nogo}

At a constant family $\Path\,A\,a\,b$, the transport component of the
structural contractum collapses by (R), and the residue is an
$\hcomp$ whose tubes are constant in the fill dimension.

\begin{lem}[Path separation]\label{lem:pathfail}
  In the algorithmic equality, over the generic context
  $(A{:}\,\U,\ a\,b{:}\,A,\ p{:}\,\Path\,A\,a\,b)$,
  \[
    \langle j \rangle\, \hcomp^{\,\iota} A\,
      [\,j{=}0 \mapsto a,\ j{=}1 \mapsto b\,]\,(p\,j)
    \;\not\aconv\; p.
  \]
  (Negative probe in \texttt{LibSwitchover.lean}; a control probe
  confirms the left-hand side is well-typed, so the failure is one of
  conversion, not typing.)  For the \emph{minimal core} the schema
  version --- $A$ any fixed core type --- is moreover proved on
  paper, unconditionally, in
  Proposition~\ref{prop:papersep}; the machine probe remains the
  witness for the full implemented calculus, where $A$ is a
  universe variable.
\end{lem}

\noindent The derivation of that equation in any equational extension
is short and fully formal.  Write $P := \Path\,A\,a\,b$ (so
$i \notin \FV(P)$) and fix the typing premises
$\Gamma \vdash A : \U$, $\Gamma \vdash a, b : A$,
$\Gamma \vdash p : P$; note $\Gamma, i \vdash P\ \mathsf{type}$ and
$\Gamma \vdash \transp^i P\,\io\,p : P$.
\begin{align}
  \Gamma &\vdash \transp^i P\,\io\,p \conv p
    && \text{by (R), $P$ $i$-free} \label{eq:d1}\\
  \Gamma &\vdash \transp^i P\,\io\,p \conv
    \langle j \rangle\, \vcomp^i A\,[\partial j \mapsto a, b]\,(p\,j)
    && \text{structural $\Path$ rule} \label{eq:d2}\\
  \Gamma &\vdash \langle j \rangle\, \vcomp^i A\,[\cdots]\,(p\,j)
    \conv \langle j \rangle\, \hcomp^{\,\iota} A\,[\cdots]\,(p\,j)
    && \text{$\vcomp = \hcomp \circ \transp$; (R); congr.}
    \label{eq:d3}\\
  \Gamma &\vdash p \conv
    \langle j \rangle\, \hcomp^{\,\iota} A\,
      [\,j{=}0 \mapsto a,\ j{=}1 \mapsto b\,]\,(p\,j)
    && \text{sym., trans.\ of \eqref{eq:d1}--\eqref{eq:d3}}
    \tag{$\dagger$}\label{eq:dagger}
\end{align}
Each step preserves the stated typing: \eqref{eq:d1} by
Lemma~\ref{lem:preservation}; \eqref{eq:d2} is the structural rule at
its stated type; in \eqref{eq:d3} the transported tubes and base are
along $i$-free lines, so (R) applies under the binder and the
congruence rules for $\hcomp$ systems, and the boundary conditions
($j{=}0$: base $p\,0 \conv a$ matches the tube; $j{=}1$
symmetrically) are exactly those of the structural rule.

Four remarks on scope, answering natural objections.  (i) The rule
used in \eqref{eq:d2} is the $\Path$-instance of the transport
unfolding common to CCHM~\cite[\S3]{CCHM} and
ABCFHL~\cite{ABCFHL}; the heterogeneous $\mathsf{PathP}$ instance
differs only in that the composition line is $A(i,j)$, and at a
constant family degenerates to the same residue.  (ii) The residue
survives in the Cartesian calculus unchanged --- no connections are
used anywhere in
\eqref{eq:d1}--\eqref{eq:dagger}.  (iii) $\eta_{\Path}$ cannot
discharge \eqref{eq:dagger}: $\eta$ cancels a path abstraction
against a path application, and the right-hand side's body is an
$\hcomp$ value, not an application of $p$.  (iv) ``Constant tubes''
means constancy in the \emph{fill} dimension $\iota$; the tubes do
depend on $j$.

\begin{thm}[No-go]\label{thm:nogo}
  \begin{enumerate}
  \item Any equational theory containing the structural
    path-transport rule, rule (R), congruence and transitivity
    derives \eqref{eq:dagger}.  Consequently the conversion checker
    of Section~\ref{sec:oper}, which refutes \eqref{eq:dagger}
    (Lemma~\ref{lem:pathfail}), is incomplete for every such theory.
  \item Equation \eqref{eq:dagger} is a constant-system
    $\hcomp$-regularity principle.  We do not claim it is literally
    an instance of a published counterexample; it is closely related
    to the regularity principles excluded by the results of Swan
    \cite{SwanRegularity} and Sattler \cite{SattlerNote}, and we
    conjecture that it fails in the standard cubical-sets semantics.
    Constructing a countermodel is stated as an open problem
    (Section~\ref{sec:conclusion}); no result of this paper depends
    on it.
  \item Restricting the structural transport rules to judgmentally
    non-constant families does not define a substitution-closed rule
    set (Section~\ref{sec:rule}); closing the restricted theory under
    substitution re-imports \eqref{eq:dagger}.  An explicit witness:
    over $\Gamma = (A : \U,\ a\,b : A,\ p : \Path\,A\,a\,b)$ the
    family $C(i) := \Path\,A\,a\,(p\,i)$ is judgmentally
    \emph{non}-constant (its normal form mentions $i$ through the
    neutral $p$), so the restricted theory applies the structural
    rule:
    $\transp^i C\,\io\,q \conv
    \langle j \rangle\, \vcomp^i A\,[\ldots]\,(q\,j)$.
    Substituting $\sigma := [\refl_a / p]$ (and $a$ for $b$) sends
    $C$ to the \emph{constant} family $\Path\,A\,a\,a$; the
    substituted instance of the derived equation must still hold in
    a substitution-closed theory, while (R) now collapses its
    left-hand side --- transitivity yields exactly
    \eqref{eq:dagger} at $\Path\,A\,a\,a$.
  \item The conflict is independent of $\Glue$: the derivation uses
    path types only.
  \end{enumerate}
\end{thm}

\noindent Theorem~\ref{thm:nogo} is, we believe, the precise content
of the folklore judgment that evaluation-time constancy detection
``does not interact well with the rest of cubical type theory'': the
obstruction is not the detection, and not specifically $\Glue$, but
that any equational marriage of (R) with the $\hcomp$-generating
structural rules collapses onto regularity for $\hcomp$.  It also
fixes the design space: the extension must be operational, with (R)
tried \emph{before} structural dispatch.

\subsection{The $\Glue$ residue}
\label{sec:glueresidue}

The same computation for $\Glue$ shows its switchover residue is the
path residue plus fiber corrections.  Let
$G := \Glue\,[\varphi \mapsto (T, w)]\,B$ with $\varphi, T, w, B$ all
$i$-free, and $u_0 : G$.  Unfold $\mathsf{transpGlue}$
\cite[\S6.2]{CCHM} at the constant line, collapsing every inner
transport by (R): the branch transport $\transp^i T\,u_0$ collapses;
the branch filler line becomes the \emph{constant} tube
$\iota \mapsto w.1\,u_0$; hence the corrected base is
\[
  a_1' \;=\; \hcomp^{\,\iota} B\,[\,\delta \mapsto w.1\,u_0\,]\,
  (\unglue\,u_0), \qquad \delta = \forall i.\,\varphi = \varphi,
\]
a constant-tube $\hcomp$, stuck off $\delta$ --- the residue of
Lemma~\ref{lem:pathfail} --- after which the fiber-point correction
adds further $\hcomp$s on the fiber $\Sigma$-type before reassembly
with $\glue$.

\section{The semantic domain, and two value equivalences}
\label{sec:domain}

This section and the next two supply the metatheory that entitles us
to call an algorithmic equality a \emph{judgmental} equality.  We
work with the concrete NbE domain of the implementation; all
definitions are stated so as to be implementation-independent.

\subsection{The domain}
\label{sec:rawdomain}

\emph{Levels} $\ell \in \mathbb{N}$ name free variables (term and
dimension alike); a \emph{depth} $d$ records the number of enclosing
binders, with the invariant that every free level of a value produced
at depth $d$ is $< d$.  \emph{Interval values} are elements of the
free De~Morgan algebra over levels, with a canonical antichain-DNF.
\emph{Values} are weak-head normal: a head constructor with value,
interval, cofibration, system and \emph{closure} arguments.
\emph{Closures} are defunctionalized: either
$\mathsf{mk}(\rho, t)$ --- an environment (a list of values) and a
term body --- or one of finitely many closure combinators
(reparametrizations, the structural-transport auxiliaries, \dots);
\emph{instantiation} $\capp_d(c, w)$ dispatches on the combinator
and, in the $\mathsf{mk}$ case, calls
$\evalop(t, \rho{+}[w])$ --- the one place evaluation re-enters, and
the reason no syntactic termination measure exists.  \emph{Neutrals}
are elimination spines on a variable level.  The \emph{dimension
action} $v\langle\xi\rangle$ of a partial map
$\xi : \mathrm{Level} \rightharpoonup \I$ applies $\xi$ to interval
parts, preserving heads and acting pointwise on closure
environments.

\subsection{Representation equivalence}
\label{sec:repeq}

In a defunctionalized NbE, the values $\evalop(A\sigma, \rho)$ and
$\evalop(A, \evalop(\sigma,\rho))$ are \emph{not} identical: the
first carries the substituted body with a shorter environment, the
second the original body with a longer one.  Every ``equal values''
statement below is therefore made with respect to:

\begin{defi}[Representation equivalence $\repeq$]\label{def:repeq}
  The largest relation such that $v \repeq v'$ implies: same head
  constructor; interval and cofibration arguments \emph{literally
  equal} (in DNF); value arguments $\repeq$; neutral spines pointwise
  (same head level, arguments $\repeq$); and for closures,
  $c \repeq c'$ iff $\capp_d(c, w) \repeq \capp_d(c', w)$ for all
  $d$ and all arguments $w$.
\end{defi}

\noindent Literal equality of interval parts is meaningful because
fresh levels are generated deterministically by depth, and the two
evaluations traverse the same term structure at the same depths
(the depth invariant of Section~\ref{sec:rawdomain}; this is made a
standing lemma in Appendix~\ref{app:cases}).

\begin{lem}[Read-back invariance]\label{lem:readback}
  $v \repeq v' \implies \quo_d(v) = \quo_d(v')$.
\end{lem}
\begin{proof}
  $\quo$ dispatches on heads, copies interval parts, and opens
  closures at generic levels --- exactly the observations $\repeq$
  preserves; coinduction over the read-back.
\end{proof}

\noindent $\repeq$ is preserved by all semantic operations
($\capp$, application, projections, path application, $\transp$,
$\hcomp$, $\glue$/$\unglue$): each dispatches on heads, interval
parts and instantiation behavior only.  The per-operation
compatibility lemmas are stated in Appendix~\ref{app:cases}.

\subsection{The Kripke logical relation}
\label{sec:valeq}

\begin{defi}[$\valeq$]\label{def:valeq}
  By induction on type values (stratified by universe level), a
  Kripke partial equivalence relation.  In this section, for
  readability, worlds are depths; the full world structure ---
  depths together with cofibration assumptions, with dimension
  substitutions and face restrictions as maps --- is what the
  clauses quantify over wherever systems and faces appear, and is
  made fully explicit for the core in
  Appendix~\ref{app:coreworlds} (for the general calculus it is
  part of component (B), as in \cite{HuberNorm}, whose
  computability predicates are likewise face-indexed families):
  \begin{itemize}
  \item $\valeq_{\mathbb{N}}$: forced to the same numeral, or both
    neutral with equal read-back;
  \item $f \valeq_{\Pi(D,C)} f'$ iff for all $d' \ge d$ and all
    $w \valeq_D w'$:
    $f\,w \valeq_{\capp(C,w)} f'\,w'$;
  \item $\Sigma$: componentwise; $\mathsf{PathP}(F,l,r)$: at every
    interval value, in the corresponding fiber;
  \item $\Glue$: the unglued parts related in the base, and on each
    live face the branch parts related in the branch type;
  \item $\U$: both codes decode to related PERs (the universe
    induction);
  \item HITs: constructor-wise (value arguments related, interval
    arguments literally equal), closed under $\hcomp$ and neutrals.
  \end{itemize}
\end{defi}

\begin{lem}\label{lem:valeqbasics}
  (i) $\repeq \subseteq \valeq_T$ for every $T$.
  (ii) $\valeq$ is a PER, invariant under $\valeq$-replacement of the
  type value.
  (iii) \emph{Conversion correctness}: $\convop_d(v, v', T)$ accepts
  iff $v \valeq_T v'$.  Soundness of $\convop$ (accepted implies
  related) is by induction on the run: each clause checks exactly a
  defining clause of $\valeq$.  Completeness (related implies
  accepted) is the standard NbE read-back argument, part of the base
  component (B) for base types and extended in
  Appendix~\ref{app:cases} for the new clause.
  In particular symmetry and transitivity of the algorithmic
  equality are inherited from the PER structure.
\end{lem}

\begin{rem}[Where clause (iii) is proved, per layer]
  \label{rem:whichlayer}
  In the \emph{relative} full calculus, clause (iii) --- in
  particular its term-level soundness direction --- is part of the
  base component (B) of Section~\ref{sec:assumptions}, inherited in
  the form stated.  For the \emph{unconditional core}, Appendix~F
  does \emph{not} re-prove term-level soundness and does not need
  it: it replaces clause (iii) by the read-back characterization of
  conversion (Lemma~\ref{lem:convrb}), type-level soundness
  (Lemmas~\ref{lem:typeidem}--\ref{lem:typeconvsound}) and
  completeness (Lemma~\ref{lem:coreconvcomp}); see
  Remark~\ref{rem:nosemconv} for why this weaker package suffices
  for every use the unconditional results make.  The two proof
  routes coexist without conflict, but they are different routes,
  and each theorem below cites the one appropriate to its layer.
\end{rem}

\section{The constancy check: specification, soundness,
  implementation}
\label{sec:check}

\subsection{Specification}

The system is \emph{defined} with the following
representation-independent check.  Let $F$ be the family value
instantiated at a fresh dimension level $\ell$ at depth $d{+}1$.

\begin{defi}[Check specification]\label{def:spec}
  $\constspec(F, \ell) :\iff \ell \notin \FV(\quo_{d+1}(F))$.
\end{defi}

\noindent By Lemma~\ref{lem:readback} the specification is invariant
under $\repeq$ --- the property that makes the substitution lemma
below go through, and that no direct traversal of representations
can have.

\subsection{Semantic soundness}

\begin{thm}[Constancy-check soundness]\label{thm:checksound}
  If $\constspec(F, \ell)$ then for all interval values $r, s$:
  $F\langle r/\ell\rangle \valeq F\langle s/\ell\rangle$ (as type
  values: they decode to the same PER).
\end{thm}
\begin{proof}
  (1) \emph{Read-back/action commutation}:
  $\quo(v\langle\xi\rangle) = \quo(v)\langle\xi\rangle$ --- the
  action touches only interval parts and closure environments, and
  read-back opens closures at generic levels disjoint from
  $\mathrm{dom}\,\xi$; induction over the read-back.
  (2) Hence with $\ell \notin \FV(\quo F)$:
  $\quo(F\langle r/\ell\rangle) = \quo(F)\langle r/\ell\rangle =
  \quo(F) = \quo(F\langle s/\ell\rangle)$.
  (3) \emph{NbE idempotence at type values}: the family value and
  the evaluation of its read-back determine the same relation.  For
  the general calculus this is part of the base component (B)
  (Appendix~\ref{app:cases}); for the minimal core it is proved as
  an independent lemma (Lemma~\ref{lem:typeidem}), by structural
  induction on the family's syntactic type --- read-back equality
  and quote--eval idempotence are distinct theorems, and we use the
  latter, at type values only.  Then
  $F\langle r/\ell\rangle \valeq \evalop(\quo F) \valeq
  F\langle s/\ell\rangle$.
\end{proof}

\noindent Theorem~\ref{thm:checksound} is what entitles the positive
branch of the evaluator to return its input: the two endpoint type
values have \emph{equal read-back} and decode to the same PER, so no
coercion is needed --- see the transport case of
Theorem~\ref{thm:fundamental}.

\subsection{The implementation realizes the specification exactly}
\label{sec:implcheck}

The kernel's occurs-check (\texttt{usesLvl}) mirrors the read-back
traversal clause by clause: the same face forcing at the head, the
same structural recursion on value and interval arguments, binder
closures instantiated at generic levels exactly as read-back's
binder cases, and the tubes of neutral $\hcomp$s likewise
instantiated, exactly as read-back's system case.  Its two
optimizations are semantically transparent: an $O(1)$
\emph{level-bound cut} --- every value carries an upper bound on its
free levels, maintained by the smart constructors; the bound's
correctness (Lemma~L1, Appendix~\ref{app:depth}) makes the cut
\emph{exact} --- and memoization.

\begin{thm}[Checker exactness, Kleene form]\label{thm:exact}
  For every value $F$ and level $\ell$: if either of
  $\texttt{usesLvl}(\ell, F)$, $\quo(F)$ is defined then both are,
  and then
  \[
    \texttt{usesLvl}(\ell, F) = \mathsf{true}
    \iff \ell \in \FV(\quo(F)).
  \]
  Hence the implemented check decides $\constspec$ exactly on
  exactly the inputs where the specification is defined; in
  particular it is invariant under $\repeq$
  (by Lemma~\ref{lem:readback}), although its traversal never
  constructs the read-back term.
\end{thm}
\begin{proof}
  Simultaneous induction on the defining big-step derivations
  (Section~\ref{sec:partiality}) of \texttt{usesLvl} and $\quo$,
  whose recursion structures are mirrored clause by clause; the full
  clause table, covering every value, neutral and
  closure-instantiation case, together with the definedness
  transfer, the completeness of the level-bound cut, and the
  transparency of memoization, is Appendix~\ref{app:exact}.
\end{proof}

\begin{rem}[An instructive archaeology]\label{rem:archaeology}
  An earlier version of the checker walked the tubes of neutral
  $\hcomp$s \emph{structurally}, over their closure environments,
  without instantiation --- a representation-sensitive
  over-approximation.  That design would have opened a real gap:
  a family the specification accepts but the implementation
  declines, making the specification take the (R)-branch and the
  implementation the structural branch --- whose results, at path
  types, are \emph{deliberately} non-convertible
  (Lemma~\ref{lem:pathfail}); one-way soundness of the check would
  then not make the kernel a realization of the specified system at
  all.  Inspection during the present revision showed that code path
  had become unreachable --- the live path already instantiates,
  mirroring read-back --- and it has been removed from the artifact,
  with the exact-mirroring property now asserted in the kernel's
  documentation.  We record this because it sharpens
  Remark~\ref{rem:whyspec} into a slogan: \emph{any deviation of the
  checker from read-back is not a harmless approximation but a
  change of judgmental equality.}
\end{rem}

\noindent Consequently the implemented evaluator \emph{is} the
specified system: every branch decision agrees
(Theorem~\ref{thm:exact}), all theorems of
Sections~\ref{sec:oper}--\ref{sec:canonicity} apply to the
implementation verbatim, and every machine probe --- positive and
negative --- is a claim about the one algorithmic equality.

\section{The operational system and its admissibility}
\label{sec:oper}

\subsection{Partiality: what kind of objects the definitions are}
\label{sec:partiality}

The functions $\evalop$, $\capp$, $\quo$, $\constspec$ (which calls
$\quo$) and $\convop$ form \emph{one} mutually recursive family of
\emph{partial} functions.  Mathematically, we define their
\emph{graphs} inductively, as big-step derivations: a derivation
certifies ``this call is defined, with this result''; determinism is
by construction of the rules, and the implementation's recursive
definitions compute exactly these graphs.  Two call cycles are
inherent and intended: instantiation re-enters evaluation on stored
closure bodies (Section~\ref{sec:rawdomain}), and the transport
clause calls the check, which reads back the family value, which
re-instantiates closures ---
$\evalop \to \constspec \to \quo \to \capp \to \evalop$.  Neither is
vicious: inductive definitions of graphs need no termination
argument; \emph{totality on well-typed inputs} is a theorem
(Corollary~\ref{cor:termination}), proved by the logical relation
and never assumed.  Three consequences for the statements below:
the closure clause of $\repeq$ (Definition~\ref{def:repeq}) is read
Kleene-style --- for all $w$, $\capp(c,w)$ and $\capp(c',w)$ are both
undefined or both defined and related; the substitution lemma is
stated in Kleene form; and the compatibility lemmas of
Appendix~\ref{app:repeq} are simultaneous inductions on the defining
derivations, with coinduction on the produced values.

Formally, $\repeq$ is the greatest fixed point of an operator
$\Phi$ on relations over the partial computation graph: for a
relation $R$ on values, $\Phi(R)$ relates $v, v'$ iff they have the
same head constructor, literally equal interval and cofibration
arguments, $R$-related value arguments, pointwise-$R$ neutral
spines, and, for each closure argument pair $(c, c')$ and every $w$,
either both $\capp(c,w)$ and $\capp(c',w)$ are undefined or both are
defined and $R$-related.  $\Phi$ is monotone (it uses $R$ only
positively), so $\repeq := \nu\Phi$ exists by Knaster--Tarski;
coinductive proofs below exhibit $\Phi$-invariant relations
(post-fixed points), and their guardedness is the fact that each
appeal to the coinduction hypothesis sits under at least one value
constructor of $\Phi$'s unfolding.  The simultaneous inductions of
Appendix~\ref{app:repeq} are ordered by the sub-derivation relation
of the big-step graphs --- an evaluation derivation strictly
contains the derivations of its premises, and instantiation
derivations are premises of the evaluation derivations that invoke
them --- which is well-founded because derivations are finite trees;
no measure on values or terms is used anywhere.

\subsection{The prioritized evaluator}

$\evalop(t, \rho)$ is the base NbE evaluator with one changed clause.
On $\transp^i A\,\varphi\,u$:
\begin{enumerate}
\item $F := \capp_{d+1}(\evalop(A, \rho), \ell)$ at a fresh dimension
  level $\ell$;
\item if $\constspec(F, \ell)$ (implementation: \texttt{usesLvl}),
  return $\evalop(u, \rho)$;
\item otherwise dispatch structurally on the head of $F$ as in the
  base evaluator ($\Pi$, $\Sigma$, $\Path$, HITs, $\Glue$ via
  $\mathsf{transpGlue}$, neutral).
\end{enumerate}
The check runs on the \emph{evaluated} family, so it fires on
families whose value is constant even when their syntax is not
(machine-checked: \texttt{strictTranspValueConstD}).  There is
exactly one definition of constancy in the system --- the one the
evaluator runs.

\begin{defi}[Typed algorithmic equality]\label{def:alg}
  $\Gamma \vdash t \aconv u : A$ iff
  \[
    \convop_{|\Gamma|}\bigl(\evalop(t, \rho_\Gamma),\
    \evalop(u, \rho_\Gamma),\ \evalop(A, \rho_\Gamma)\bigr)
  \]
  accepts, where
  $\rho_\Gamma$ is the generic environment of $\Gamma$ (each variable
  a fresh neutral at its type).
\end{defi}

\begin{defi}[The system of record $T_{\mathrm{alg}}$]\label{def:talg}
  $T_{\mathrm{alg}}$ is the type system with the rules of
  Appendix~\ref{app:rules} whose conversion rule is mediated by
  $\aconv$ (Definition~\ref{def:alg}); it is the system the
  bidirectional checker implements.  It is \emph{not} presented by an
  equational theory on terms: by Theorem~\ref{thm:nogo} no
  substitution-closed rule-generated presentation containing both
  (R) and the structural rules can coincide with it.
\end{defi}

\subsection{The substitution lemma}

\begin{lem}[Substitution/environment coherence]\label{lem:subst}
  For every term $t$, substitution $\sigma$ and environment $\rho$:
  if either of $\evalop(t\sigma, \rho)$,
  $\evalop(t, \evalop(\sigma, \rho))$ is defined then both are, and
  \[
    \evalop(t\sigma, \rho) \;\repeq\; \evalop(t, \evalop(\sigma, \rho)).
  \]
\end{lem}
\begin{proof}
  Structural induction on $t$.  All cases except transport are the
  standard NbE compositionality cases: value formers produce equal
  heads with $\repeq$ arguments; the closure-forming cases are the
  closure clause of Definition~\ref{def:repeq} (instantiating either
  closure continues with the induction hypothesis); interval parts
  are literally equal by the depth-determinacy of fresh levels.

  For $t = \transp^i A\,\varphi\,u$: the two family values
  \[
    F_L := \evalop(A\sigma,\rho)[\ell], \qquad
    F_R := \evalop(A, \evalop(\sigma,\rho))[\ell]
  \]
  satisfy $F_L \repeq F_R$ (induction hypothesis for $A$ plus
  $\capp$-compatibility), with the \emph{same} fresh $\ell$ (depth
  determinacy).  The check is the specification of
  Definition~\ref{def:spec}, which is $\repeq$-invariant
  (Lemma~\ref{lem:readback}); \emph{hence the two sides take the same
  branch}.  Positive branch: induction hypothesis for $u$.  Negative
  branch: the same structural dispatch on the same head, closing by
  the induction hypotheses and the per-operation compatibility
  lemmas.
\end{proof}

\begin{rem}\label{rem:whyspec}
  The proof shows why the check \emph{must} be specified
  representation-independently: an implementation-level check
  (a traversal of representations) is not a function of the
  $\repeq$-class --- the two sides of the lemma present the same
  family in different representations, and a
  representation-sensitive check could take different branches,
  breaking the lemma at exactly the transport case.  Specification
  vs.\ implementation is not pedantry here; it is the load-bearing
  design decision.  (For the artifact no gap remains: the implemented
  checker decides the specification exactly, Theorem~\ref{thm:exact}
  --- and Remark~\ref{rem:archaeology} records how close it once came
  to having one.)
\end{rem}

\subsection{The fundamental theorem, and termination}

\begin{defi}[Semantic typing]\label{def:semtyping}
  $\Gamma \vDash t : A$ iff for all depths and all pairs of related
  $\Gamma$-environments $\rho \valeq \rho'$: $\evalop(t, \rho)$ is
  \emph{defined} (evaluation terminates) and
  $\evalop(t, \rho) \valeq_{\evalop(A,\rho)} \evalop(t, \rho')$.
\end{defi}

\begin{thm}[Fundamental theorem]\label{thm:fundamental}
  If $\Gamma \vdash t : A$ in $T_{\mathrm{alg}}$ then
  $\Gamma \vDash t : A$.
\end{thm}
\begin{proof}
  Induction on the typing derivation.  All $T_0$ rules are the base
  component (B), in environment form (the enumeration, with the
  clause-by-clause correspondence to \cite{HuberNorm}, is
  Appendix~\ref{app:cases}).  Two cases are specific to
  $T_{\mathrm{alg}}$.

  \emph{The conversion rule.}  From $\Gamma \vdash t : B$ and
  $\Gamma \vdash B \aconv A$: by IH, $\evalop(t,\rho)$ is defined and
  related at $\evalop(B,\rho)$; conversion correctness
  (Lemma~\ref{lem:valeqbasics}(iii)) turns the accepted conversion
  into $\evalop(B,\rho) \valeq \evalop(A,\rho)$, and $\valeq$ is
  invariant under related type values (Lemma~L10), closing the case.
  This is where defining the system of record with the
  \emph{algorithmic} equality is essential and not cosmetic: had the
  conversion rule used derivability in $T_{\mathrm{eq}}$, this case
  would demand invariance of the predicates along
  \eqref{eq:dagger}-type equations, which the algorithm refutes ---
  the fundamental theorem would be false as stated.

  \emph{The transport clause} of the prioritized evaluator:
  \begin{itemize}
  \item \emph{Positive branch.}  By Theorem~\ref{thm:checksound} the
    endpoint type values decode to the same PER --- indeed they have
    equal read-back --- so the induction hypothesis for $u$ gives
    exactly the required relatedness at the result type.  Note no
    closure property of the predicates under judgmental equality is
    invoked; the equality of PERs is by read-back identity.  This is
    where the specification (Definition~\ref{def:spec}) pays a
    second time.
  \item \emph{Negative branch.}  The dispatch clauses are the base
    proof's structural transport cases.  Their inner recursive
    transports (along domain lines, cod fibers, tube lines) are at
    \emph{structurally smaller semantic type values}, so they are
    discharged by the inner induction on type values by which
    $\valeq$ and the base transport case are defined --- \emph{not}
    by the outer induction on typing derivations.  The proof is thus
    a double induction: outer on the derivation, inner on type
    values; the evaluator's own recursion is mirrored by the inner
    induction, which is well-founded by the stratified definition of
    $\valeq$ (Definition~\ref{def:valeq}).  Termination of the
    dispatch is part of the conclusion, not an input.
  \end{itemize}
\end{proof}

\begin{cor}[Termination]\label{cor:termination}
  Evaluation, the constancy check, conversion and read-back
  terminate on well-typed inputs.
\end{cor}
\begin{proof}
  Termination of $\evalop$ is part of
  Definition~\ref{def:semtyping} and hence of
  Theorem~\ref{thm:fundamental}.  The check traverses a value that is
  semantically well-typed; its binder instantiations are
  instantiations of computable closures at generic arguments, which
  terminate by the same theorem; system closures are walked without
  instantiation.  $\convop$ and $\quo$ recurse along the stratified
  type-value structure.  No syntactic measure is claimed anywhere:
  this is the standard normalization-proof route
  \cite{HuberNorm,SterlingAngiuli}, and the only one available, since
  closure instantiation re-enters evaluation
  (Section~\ref{sec:rawdomain}).
\end{proof}

\subsection{Admissibility of the algorithmic equality}

\begin{thm}[Admissible conversion relation]\label{thm:admissible}
  Relative to the base component (B), the typed algorithmic equality
  of Definition~\ref{def:alg}:
  \begin{enumerate}
  \item is an equivalence relation on well-typed terms
    (reflexivity from Theorem~\ref{thm:fundamental} at the generic
    environment; symmetry and transitivity from the PER structure via
    Lemma~\ref{lem:valeqbasics}(iii));
  \item is a congruence: for every term former, relatedness of the
    immediate subterms implies relatedness of the whole (each former
    evaluates to a value former whose $\valeq$-clause is
    componentwise; the former-by-former list, including both
    transport branches, is Appendix~\ref{app:cases});
  \item is stable under weakening (Kripke monotonicity of $\valeq$
    in the depth);
  \item is stable under substitution: $\Gamma \vdash t \aconv u : A$
    and $\Delta \vdash \sigma : \Gamma$ imply
    $\Delta \vdash t\sigma \aconv u\sigma : A\sigma$
    (Lemma~\ref{lem:subst} composed with
    Lemma~\ref{lem:valeqbasics}(i,iii) and
    Theorem~\ref{thm:fundamental} applied to $\sigma$);
  \item is stable under context conversion ($\Gamma \aconv \Gamma'$
    pointwise implies the generic environments are $\valeq$-related,
    and $\valeq$ is invariant under related type values);
  \item preserves and respects typing in $T_{\mathrm{alg}}$: if
    $\Gamma \vdash t : A$ and $\Gamma \vdash A \aconv B$ then
    $\Gamma \vdash t : B$ (this \emph{is} $T_{\mathrm{alg}}$'s
    conversion rule, and the fundamental theorem shows it harmless),
    and the bidirectional checker's conversion sites decide exactly
    this relation.
  \end{enumerate}
\end{thm}

\noindent Theorem~\ref{thm:admissible} is the license to call the
algorithmic equality a \emph{judgmental} equality: it has every
structural property the conversion rule of a dependent type theory
requires.

\section{Soundness, and what completeness would mean}
\label{sec:soundness}

\begin{lem}[Step soundness]\label{lem:stepsound}
  Write $\hat\rho$ for the substitution quoting an environment
  pointwise, and $t[\hat\rho]$ for the result of applying it.  If
  $\evalop(t, \rho)$ is defined with value $v$, then
  $t[\hat\rho] \conv \quo(v)$ is derivable in
  $T_{\mathrm{eq}} = T_0 + (\mathrm{R})$; and if
  $\convop(v, v', T)$ accepts then $\quo(v) \conv \quo(v')$ is
  derivable in $T_{\mathrm{eq}}$.
\end{lem}
\begin{proof}
  Induction on the big-step derivation
  (Section~\ref{sec:partiality}), with the read-back statements of
  the sub-computations supplied by the same induction.  Each
  evaluator rule is matched by a $T_{\mathrm{eq}}$ equation:
  $\beta$-steps and projections by the $\beta$ rules; face collapses
  by the boundary rules; the structural transport dispatches by the
  corresponding structural rules of Section~\ref{sec:rules},
  instantiated at the quoted premises; $\hcomp$ steps by the
  composition rules; $\eta$-directed conversion clauses by the
  $\eta$ rules, and interval comparisons by the free De~Morgan
  axioms (sound for $T_{\mathrm{eq}}$, which contains them).

  The one step needing care is the \emph{positive transport branch}:
  $\evalop(\transp^i A\,\varphi\,u, \rho) = \evalop(u,\rho)$ when
  $\constspec(F, \ell)$ for $F$ the family value at the fresh
  $\ell$.  We must exhibit a $T_{\mathrm{eq}}$ derivation of
  $(\transp^i A\,\varphi\,u)[\hat\rho] \conv u[\hat\rho]$.  Take
  $A_0 := \quo(F)$, a term not containing (the variable
  corresponding to) $\ell$, by the definition of $\constspec$.  By
  the induction hypothesis applied to the family's evaluation and
  instantiation, $A[\hat\rho] \conv A_0$ is derivable in
  $T_{\mathrm{eq}}$ (in the context extended by $\ell$).  So
  $A[\hat\rho]$ is \emph{judgmentally $i$-constant} in the sense of
  Definition~\ref{def:const} --- the premise of rule (R) --- with
  witness $A_0$, and (R) yields
  $(\transp^i A\,\varphi\,u)[\hat\rho] \conv u[\hat\rho]$.  The
  induction hypothesis for $u$ finishes the step.
\end{proof}

\begin{thm}[Soundness]\label{thm:soundness}
  Algorithmic equality is contained in derivability in
  $T_{\mathrm{eq}}$: if $\Gamma \vdash t \aconv u : A$ then
  $\Gamma \vdash t \conv u : A$ is derivable in $T_{\mathrm{eq}}$.
\end{thm}
\begin{proof}
  By Definition~\ref{def:alg}, $\convop$ accepts the values of $t$
  and $u$ at the generic environment; Lemma~\ref{lem:stepsound}
  gives $T_{\mathrm{eq}}$-derivations of
  $t \conv \quo(\evalop\,t)$, $u \conv \quo(\evalop\,u)$
  (the generic environment quotes to the identity substitution) and
  $\quo(\evalop\,t) \conv \quo(\evalop\,u)$; compose by
  transitivity.
\end{proof}

\noindent The asymmetry between Theorem~\ref{thm:soundness} and
Remark~\ref{rem:incomplete} below is the paper's central design
fact: $T_{\mathrm{alg}}$'s equality is \emph{strictly weaker} than
$T_{\mathrm{eq}}$'s (the no-go equation separates them), yet every
algorithmic step is $T_{\mathrm{eq}}$-sound; the algorithm selects a
consistent, substitution-stable sub-theory of an equational theory
that is itself unusable.

\begin{rem}[No completeness --- necessarily]\label{rem:incomplete}
  By Theorem~\ref{thm:nogo}(1), $T_0+(\mathrm{R})$ derives
  \eqref{eq:dagger}, which $\convop$ refutes; the checker is
  incomplete for $T_0+(\mathrm{R})$, provably and by design.  The
  system of record is the algorithmic equality itself ---
  legitimately so, by Theorem~\ref{thm:admissible} ---
  and $T_0+(\mathrm{R})$ is its sound equational
  over-approximation.  We therefore do not use the phrase ``decision
  procedure'' for judgmental equality anywhere: the algorithmic
  equality is decidable by running its terminating checker
  (Corollary~\ref{cor:termination}), and is soundly included in
  $T_0+(\mathrm{R})$; nothing more is claimed.
\end{rem}

\section{Canonicity}
\label{sec:canonicity}

\begin{thm}[Canonicity]\label{thm:canonicity}
  Relative to the base component (B), in both the Cartesian and the
  De~Morgan presentations --- in the Cartesian fragment the
  judgmental-theory part of (B) follows from \cite{SterlingAngiuli},
  the evaluator-level bridge remaining open
  (Appendix~\ref{app:bridge}): every closed $T_{\mathrm{alg}}$-term
  $t : \mathbb{N}$ evaluates to a numeral; the algorithmic equality
  is consistent ($0 \not\aconv 1$); and $\J\,d\,\refl \aconv d$.
\end{thm}
\begin{proof}
  The closed instance of Theorem~\ref{thm:fundamental}: the empty
  environment is trivially related to itself, so
  $\evalop(t) \valeq_{\mathbb{N}} \evalop(t)$, and the
  $\valeq_{\mathbb{N}}$ clause forces a numeral (there are no closed
  neutrals); numeral inversion is
  Appendix~\ref{app:cases}.  Consistency: $0$ and $1$ evaluate to
  distinct numerals, and $\convop$ separates them.  For $\J$:
  substituting $p := \refl$ makes the transport family's value
  read back without the transport dimension, so
  $\constspec$ holds and the positive branch returns $d$'s value;
  machine-checked end to end (\texttt{strictJReflD}).
\end{proof}

\begin{cor}[Unconditional minimal core]\label{thm:core}
  Over the minimal fragment $\Pi, \Sigma, \Path, \mathbb{N}$ with
  $\transp$ and $\hcomp$ (no $\Glue$, no universes, no HITs),
  Theorems~\ref{thm:fundamental}, \ref{thm:admissible}
  and~\ref{thm:canonicity} hold \emph{with no base-component
  assumption}: Appendix~\ref{app:core} constructs the logical
  relation for the fragment by structural recursion on syntactic
  types and proves every (B)-marked lemma for it, former by former.
  In particular $\J\,d\,\refl \aconv d$ --- $\J$ being defined from
  paths, connections and transport alone --- is unconditional
  (Proposition~\ref{prop:J}).
\end{cor}

\subparagraph*{Inheritance inventory.}  Reused from
\cite{HuberNorm} (as component (B)): the predicate clauses for base
formers, $\hcomp$-computability (the composition clauses never
inspect transport redexes), read-back correctness, numeral
inversion.  New in this paper: representation equivalence and its
compatibility lemmas (Section~\ref{sec:repeq}); the check
specification, its soundness (Theorem~\ref{thm:checksound}) and the
checker-exactness theorem (Theorem~\ref{thm:exact}); the
substitution lemma through the transport case
(Lemma~\ref{lem:subst}); the positive transport case of the
fundamental theorem; the explicit double-induction discharge of the
negative case; and the admissibility theorem
(Theorem~\ref{thm:admissible}).

\section{The strictness catalogue}
\label{sec:catalogue}

Machine-checked excerpt (positive rows: an inline $\refl$ witness is
accepted; negative rows: the analogous witness is rejected; all rows
are claims about the one algorithmic equality --- implementation and
specification coincide by Theorem~\ref{thm:exact}):

\begin{center}
\small
\begin{tabular}{ll}
  law & holds? \\\hline
  connection reparametrizations (idempotence, involution,
    absorption, $i \wedge 0$) & yes \\
  De Morgan duality square & yes \\
  $\eta$ for $\Pi$, $\Sigma$, $\Path$ & yes \\
  $\mathsf{symm}\,\refl \conv \refl$;\
    $\mathsf{symm}\circ\mathsf{symm} \conv \mathrm{id}$;\
    $\mathsf{cong}$ laws & yes \\
  transport along constant families (syntactic and value-level; at
    $\U$; at path families) & yes \\
  switchover contracta at $\Pi$, $\Sigma$, dependent $\Sigma$,
    argument-wise HIT constructors (incl.\ $\mathsf{merid}$) & yes \\
  $\J\,d\,\refl \conv d$ & \textbf{yes} \\
  $\mathsf{transport}\,(\ua\,\mathsf{idEquiv})\,x \conv x$
    (via $\mathsf{transpGlue}$) & yes \\
  $\mathsf{trans}\,\refl\,\refl \conv \refl$
    ($\hcomp$-regularity) & no \\
  path switchover contractum \eqref{eq:dagger} & no \\
  $\mathsf{cong}$/$\mathsf{symm}$ distribution over
    $\mathsf{trans}$ & no \\
  $\langle i\rangle\, p(i \wedge \neg i) \conv \refl$
    (Appendix~\ref{app:reparam}) & no \\
\end{tabular}
\end{center}

\section{Implementation and artifact}
\label{sec:impl}

The kernel is a self-contained CCHM-style implementation of about
five thousand lines of Lean~4 (toolchain \texttt{4.31.0}),
independent of Lean's own definitional equality: the NbE core with
defunctionalized closures, bidirectional type checking, the De~Morgan
interval with the antichain-DNF decision procedure,
$\Glue$/univalence with $\mathsf{transpGlue}$, and the higher
inductive types of Section~\ref{sec:hits}.  The prioritized transport
of Section~\ref{sec:oper} is the evaluator's transport clause; the
occurs-check is implemented as described in
Section~\ref{sec:implcheck}.  A $\sim$300-definition object-language
library exercises the rule throughout ($\pi_1(S^1) \cong \mathbb{Z}$,
untruncated Eckmann--Hilton and Mac\,Lane coherence, a groupoid
classifying-space HIT with its classification theorem), validated by
a differential harness fingerprinting the normal forms of every
library definition.

\subparagraph*{Artifact and its scope.}  The repository --- kernel,
catalogue (\texttt{LibStrictness.lean}), switchover experiments
(\texttt{LibSwitchover.lean}), library, harness, and a row-by-row
paper-to-probe correspondence (\texttt{ARTIFACT.md}) --- is public at
\begin{center}
  \url{https://github.com/r-shibusawa/cubical-strict-j}
\end{center}
release \texttt{v1.0.0}, commit \texttt{6c5904b}, archived at
DOI \href{https://doi.org/10.5281/zenodo.21405962}%
{\texttt{10.5281/zenodo.21405962}}; pinned toolchain
Lean \texttt{4.31.0}; build with \texttt{lake build}: every catalogue
and switchover claim is an elaboration-time guard, so a successful
build is the verification run (continuous integration replays it on
every commit).  We distinguish explicitly: the artifact provides
\emph{machine-checked object-language witnesses} (all
algorithmic-equality claims) and \emph{executable regression tests}
(the differential harness); it does \emph{not} provide machine-checked
metatheory --- the theorems of
Sections~\ref{sec:switchover}--\ref{sec:canonicity} are pen-and-paper
results about the specified system, which the kernel implements
exactly (Theorem~\ref{thm:exact}); a kernel passing its test suite
remains evidence, not proof, of its own soundness.

\section{Related work}
\label{sec:related}

\subparagraph*{Constancy formulas.}  The $\varphi$-discipline is from
ABCFHL \cite{ABCFHL}; Cubical Agda \cite{VezzosiMortbergAbel} fixes
$\varphi$ at elaboration time, whence the propositional
\textsf{transportRefl}.  Rule (R) is the conversion-closure of
$(\mathrm{S}_{\ii})$; Remark~\ref{rem:explicit} locates exactly when
this is a proper extension.

\subparagraph*{Regularity.}  CCHM \cite{CCHM} initially claimed
regular composition; the universe error was found by Sattler
\cite{SattlerNote}, and Swan \cite{SwanRegularity} separated path and
identity types in presheaf models.  Those results concern
$\hcomp$/composition; Theorem~\ref{thm:nogo} connects them to
transport-level constancy inference: prioritization is exactly what
prevents the equational collapse onto $\hcomp$-regularity.  Recent
work towards computational UIP in Cubical Agda \cite{UIPCubical}
documents the practical cost of missing regularity.

\subparagraph*{Systems with regular coercion.}  Arend \cite{Arend}
has a regular \textsf{coe} over a non-cubical interval; the folklore
obstruction to porting this to full cubical theories --- which
Theorem~\ref{thm:nogo} makes precise --- was the interaction with
structural transport and $\Glue$.  XTT and observational theories
\cite{StickyXTT} obtain a strict $\J$ by strengthening equality
globally, at the cost of the univalent universe; the present system
keeps the ABCFHL setting unchanged.

\subparagraph*{Normalization and NbE metatheory.}  The logical
relation of Section~\ref{sec:domain} is the standard Kripke PER
structure of NbE correctness proofs, transposed to the cubical
domain; the environment-based fundamental theorem follows Huber
\cite{HuberNorm}, and the Cartesian discharge of the base component
is Sterling--Angiuli \cite{SterlingAngiuli} (normal-form
specification in \cite{HuangNF}).  What is new is confined to the
transport clause: the representation-independent check specification,
its soundness theorem, and the resulting substitution and
admissibility results.

\section{Conclusion and open problems}
\label{sec:conclusion}

Evaluation-time constancy inference is an old idea.  The analysis
here shows why its naive equational formulations could not work ---
they derive constant-system $\hcomp$-regularity, at path types
already (Theorem~\ref{thm:nogo}) --- and that the operational
formulation, done with a representation-independent check
specification, supports the full structural metatheory expected of a
judgmental equality (Theorem~\ref{thm:admissible}) together with
termination and canonicity (Theorem~\ref{thm:canonicity}), and
computes $\J$ on $\refl$ --- \emph{unconditionally on the minimal
core} $\Pi, \Sigma, \Path, \mathbb{N}$
(Theorem~\ref{thm:core}, Proposition~\ref{prop:J}), and
\emph{relative to the delimited base component} for the full
implemented univalent calculus.  We do not claim an unconditional
strict $\J$ for the full univalent system; the three-layer division
of Section~\ref{sec:intro} is the paper's claim, exactly.

Open problems, in decreasing order of importance.
(1)~\emph{The base component beyond the core}: for the core fragment
of Appendix~\ref{app:core} the base component is discharged in this
paper and the positive results are unconditional; beyond it, the
judgmental-theory part of the Cartesian base component follows from
\cite{SterlingAngiuli}, while the concrete evaluator-level bridge
(Appendix~\ref{app:bridge}), universes, $\Glue$ and the wider HIT
roster remain to be discharged --- on the De~Morgan side, in full.
(2)~\emph{A semantic countermodel for \eqref{eq:dagger}}: upgrading
Theorem~\ref{thm:nogo}(2) from ``closely related to excluded
principles'' to a refutation in a concrete cubical-sets model.
(3)~\emph{Model-theoretic status of the positive system}: whether a
cubical-sets-style model validates the specified algorithmic
equality.
(4)~\emph{Characterization}: is there any rule-generated presentation
of the operational system?  We conjecture not.
(5)~\emph{Boundary-corrected HIT signatures}: extending the
switchover classification of Section~\ref{sec:hits} beyond the
argument-wise class.

\subparagraph*{Acknowledgements.}  Three rounds of critical review of
earlier drafts materially reshaped this paper: the observation that
non-constancy is not substitution-stable prompted the switchover
experiments that became Theorem~\ref{thm:nogo}, and the demand for
the structural metatheory of the algorithmic equality prompted
Sections~\ref{sec:domain}--\ref{sec:canonicity} in their present
form.

\appendix

\section{The base calculus in full}
\label{app:rules}

This appendix fixes $T_0$ completely: grammar, judgments, typing,
and the equational rules, matching the artifact's kernel former by
former.  ($\Gamma \vdash \varphi\ \mathsf{cof}$ premises and routine
congruence rules are elided from displays.)

\subsection{Grammar and judgments}

\[
\begin{array}{r@{\ }l}
  r, s ::=& i \mid 0 \mid 1 \mid r \wedge s \mid r \vee s \mid \neg r
  \qquad\qquad
  \varphi, \psi ::= (r{=}0) \mid (r{=}1) \mid \varphi \wedge \psi
    \mid \varphi \vee \psi \mid \io \mid \ii \\[2pt]
  A, B, t, u ::=& x \mid \U_k
    \mid \Pi x{:}A.B \mid \lambda x.t \mid t\,u
    \mid \Sigma x{:}A.B \mid (t, u) \mid \fst t \mid \snd t \\
  &\mid \mathsf{PathP}\,(\langle i \rangle A)\,t\,u
    \mid \langle i \rangle t \mid t\,r
    \mid \mathbb{N} \mid \mathsf{ze} \mid \mathsf{su}\,t
    \mid \mathsf{natrec} \\
  &\mid \transp^i A\,\varphi\,u
    \mid \hcomp^{\,i} A\,[\overrightarrow{\varphi \mapsto t}]\,u
    \mid \Glue\,[\overrightarrow{\varphi \mapsto (T, w)}]\,B \\
  &\mid \glue \mid \unglue
    \mid \text{HIT formers (Section~\ref{sec:hits})}
\end{array}
\]
Judgments: $\Gamma\ \mathsf{ctx}$;\ \ $\Gamma \vdash A\ \mathsf{type}_k$
(equivalently $\Gamma \vdash A : \U_k$, universes \`a la Russell,
cumulative);\ \ $\Gamma \vdash t : A$;\ \ $\Gamma \vdash
\varphi\ \mathsf{cof}$; and the \emph{algorithmic} equality
$\Gamma \vdash t \aconv u : A$ of Definition~\ref{def:alg}, used in
the conversion rule.  Contexts extend by term variables, interval
variables and cofibration restrictions.  Interval expressions are
identified up to free De~Morgan equality (decided by antichain-DNF);
the implementation additionally provides
$\mathsf{Int}, \mathsf{Unit}, \mathsf{Empty}$, binary sums and lists,
whose rules are routine copies of the $\mathbb{N}$ pattern.

\subsection{Structural core}

\[
  \frac{\Gamma, x{:}A \vdash t : B}
       {\Gamma \vdash \lambda x.t : \Pi x{:}A.B}
  \qquad
  \frac{\Gamma \vdash t : \Pi x{:}A.B \quad \Gamma \vdash u : A}
       {\Gamma \vdash t\,u : B[u/x]}
  \qquad
  (\beta,\ \eta_\Pi:\ \lambda x.\,t\,x \conv t)
\]
\[
  \frac{\Gamma \vdash t : A \quad \Gamma \vdash u : B[t/x]}
       {\Gamma \vdash (t,u) : \Sigma x{:}A.B}
  \qquad
  (\fst/\snd\ \text{projections};\ \beta;\
   \eta_\Sigma:\ (\fst t, \snd t) \conv t)
\]
\[
  \frac{\Gamma, i{:}\I \vdash t : A}
       {\Gamma \vdash \langle i \rangle t :
         \mathsf{PathP}\,(\langle i \rangle A)\,t(\io)\,t(\ii)}
  \qquad
  \frac{\Gamma \vdash p : \mathsf{PathP}\,(\langle i \rangle A)\,l\,r
        \quad \Gamma \vdash s : \I}
       {\Gamma \vdash p\,s : A(s/i)}
\]
\[
  (\beta;\ p\,\io \conv l,\ p\,\ii \conv r;\
   \eta_{\Path}:\ \langle i \rangle\, p\,i \conv p)
\]
\[
  \frac{}{\Gamma \vdash \mathbb{N} : \U_0}
  \qquad
  \mathsf{natrec}\ \text{with its two computation rules}
  \qquad
  \frac{\Gamma \vdash A : \U_k}{\Gamma \vdash A : \U_{k+1}}
\]

\subsection{Kan operations}

Transport was displayed in Section~\ref{sec:rules}: the typing rule
with explicit $\varphi$, the collapse rule $(\mathrm{S}_{\ii})$, and
the structural rules for $\Pi$, $\Sigma$, $\Path$.  The remaining
structural transport rules: at $\U_k$, $\mathbb{N}$ and base types,
$\transp^i A\,\varphi\,u \conv u$ (the line is constant by
typing); at HIT formers, argument-wise per
Section~\ref{sec:hits}; at $\Glue$, the $\mathsf{transpGlue}$
unfolding of \cite[\S6.2]{CCHM} (its constant-family degeneration is
computed in Section~\ref{sec:glueresidue}); on neutral lines,
transport is neutral.

Homogeneous composition:
\[
  \frac{\Gamma \vdash A\ \mathsf{type} \qquad
        \Gamma \vdash u : A \qquad
        \Gamma, \varphi_k, i{:}\I \vdash t_k : A\ \ (\forall k)}
       {\Gamma \vdash \hcomp^{\,i} A\,
         [\overrightarrow{\varphi_k \mapsto t_k}]\,u : A}
\]
with side conditions
$\Gamma, \varphi_k \vdash t_k(\io) \conv u$ and
$\Gamma, \varphi_k \wedge \varphi_l, i \vdash t_k \conv t_l$,
with $\Gamma, \varphi_k \vdash
\hcomp\,[\cdots]\,u \conv t_k(\ii)$, and
$\hcomp\,[\ii \mapsto t]\,u \conv t(\ii)$; structural rules push
$\hcomp$ into $\Pi$/$\Sigma$/$\Path$ components, compute at
$\mathbb{N}$ and constructors, and are constructors at HITs
($\hcomp$-cells).  \emph{No} constancy rule for $\hcomp$ exists
(Theorem~\ref{thm:nogo} is why none may be added).
$\vcomp^i A\,[\vec\varphi \mapsto \vec t]\,u :=
\hcomp^{\,i} A(\ii)\,[\vec\varphi \mapsto
\transp\text{-filled } \vec t]\,(\transp^i A\,\io\,u)$.

\subsection{$\Glue$}

\[
  \frac{\Gamma \vdash B : \U_k \quad
        \Gamma, \varphi \vdash T : \U_k \quad
        \Gamma, \varphi \vdash w : \mathsf{Equiv}\,T\,B}
       {\Gamma \vdash \Glue\,[\varphi \mapsto (T,w)]\,B : \U_k}
  \qquad
  \Gamma, \varphi \vdash \Glue\,[\varphi \mapsto (T,w)]\,B \conv T
\]
with $\glue$/$\unglue$ introduction and elimination, their $\beta$
rule, the face collapses, and univalence realized as $\ua$ by gluing
along an equivalence; $\mathsf{Equiv}$ is the standard
contractible-fibers definition (no universe needed, as in
\cite{SterlingAngiuli}).

\subsection{Conversion rule}

\[
  \frac{\Gamma \vdash t : A \qquad \Gamma \vdash A \aconv B}
       {\Gamma \vdash t : B}
\]
The equality in the conversion rule is the algorithmic equality ---
legitimately, by Theorem~\ref{thm:admissible}; the bidirectional
checker of the artifact implements exactly this rule at its
switching points.

\section{Reparametrization coherence}
\label{app:reparam}

A self-contained result delimiting the strict/weak boundary from
below: reparametrizations are strict, and whatever generates an
$\hcomp$ is not --- as Theorem~\ref{thm:nogo} explains, this is no
accident.  Fix the generic context
\[
  \Gamma = (A : \U,\ a\,b : A,\ p : \Path\,A\,a\,b).
\]

\begin{thm}[Reparametrization coherence]\label{thm:reparam}
  Let $\varphi, \psi \in \DM{i_1,\dots,i_n}$ agree at every corner
  $\varepsilon \in \{0,1\}^n$.  Then
  \[
    \langle i_1 \cdots i_n \rangle\, p(\varphi)
    \;\aconv\;
    \langle i_1 \cdots i_n \rangle\, p(\psi)
    \qquad\Longleftrightarrow\qquad
    \varphi = \psi \ \text{in}\ \DM{i_1,\dots,i_n}.
  \]
\end{thm}
\begin{proof}
  ($\Leftarrow$) Conversion instantiates binders at generic
  dimensions $\vec\ell$; $i_k \mapsto \ell_k$ is a De~Morgan
  homomorphism, and neutral path applications compare interval
  arguments by canonical antichain-DNF, sound and complete for free
  De~Morgan equality; if the common value is $0/1$, both sides
  collapse to the same endpoint.  ($\Rightarrow$) In the generic
  context $p, a, b$ are distinct neutrals.  If neither side
  collapses, DNF equality of $\varphi(\vec\ell)$ and
  $\psi(\vec\ell)$ is demanded, and generators map to generators, so
  $\varphi = \psi$.  If exactly one side collapses, a variable is
  compared against a neutral application and conversion fails.  If
  both collapse, the corner hypothesis forces the same endpoint.
\end{proof}

\noindent In particular
$\langle i \rangle\, p(i \wedge \neg i) \not\aconv \refl_a$: the
absence of excluded middle in the free De~Morgan algebra is reflected
in the algorithmic equality (machine-checked, both directions).

\section{Lemmas: statements and proofs}
\label{app:cases}

This appendix carries out the routine inductions behind
Sections~\ref{sec:domain}--\ref{sec:canonicity}.  A structural remark
organizes everything:

\begin{rem}[The domain is unchanged]\label{rem:samedomain}
  Rule (R) alters the \emph{evaluator}, not the \emph{domain}: it
  prunes transport redexes and introduces no new value forms.  The
  values, neutrals, closures, read-back and conversion of the
  extended system are literally those of the base system.
  Consequently every base lemma \emph{about values} --- the PER laws,
  Kripke monotonicity, conversion correctness, read-back correctness,
  numeral inversion --- transfers verbatim; the extension is audited
  entirely in (i) the evaluator clauses (the fundamental theorem,
  Appendix~\ref{app:fund}) and (ii) the check
  (Section~\ref{sec:check}).  Lemmas below marked (B) are of the
  first kind and constitute the base component of
  Section~\ref{sec:assumptions}; lemmas marked (new) are of the
  second kind and are proved here in full.
\end{rem}

\subsection{Depth and freshness}
\label{app:depth}

\begin{lem}[L1: depth invariant (new)]\label{lem:L1}
  Every free level of a value produced by evaluation at depth $d$ is
  $< d$; the level bound $\mathrm{lb}(v)$ maintained by the smart
  constructors satisfies $\FV(v) \subseteq [0, \mathrm{lb}(v))$ and
  $\mathrm{lb}(v) \le d$.
\end{lem}
\begin{proof}
  Induction on the evaluation step producing $v$.  Environment
  entries were produced at outer depths, so the invariant holds for
  them by hypothesis; a fresh level is allocated as the current depth
  $d$ and only inside a binder, i.e.\ at inner depth $d{+}1$, where
  $d < d{+}1$; every smart constructor sets
  $\mathrm{lb} := \max$ of its parts' bounds, and no operation
  manufactures a level except fresh allocation.  The level-bound cut
  of Section~\ref{sec:implcheck} ($\ell \ge \mathrm{lb}(v) \implies
  \ell \notin \FV(v)$) is immediate.
\end{proof}

\begin{lem}[L2: depth determinacy (new)]\label{lem:L2}
  Let two evaluation runs traverse the same term spine at the same
  depth (in possibly different environments).  Then they allocate the
  same fresh levels at the same positions.
\end{lem}
\begin{proof}
  Fresh levels are, by construction, \emph{equal to the depth at
  allocation}; depth evolves only by crossing a binder of the term
  spine being traversed, which is the same on both runs.  (This is
  the reason Definition~\ref{def:repeq} may demand \emph{literal}
  equality of interval parts: the interval values produced by two
  $\repeq$-related computations are built from the same generic
  levels by the same De~Morgan operations.)
\end{proof}

\subsection{Compatibility of representation equivalence}
\label{app:repeq}

All compatibility lemmas are proved by one simultaneous coinduction:
the relation
\[
  \mathcal{R} \;:=\; \{\,(\mathsf{op}(\vec v),\ \mathsf{op}(\vec v'))
    \mid \mathsf{op} \text{ a semantic operation},\ \vec v \repeq
    \vec v'\,\}
\]
is shown to be a $\repeq$-bisimulation, by case analysis on
$\mathsf{op}$ and on the (common) head of the principal arguments.
We record the cases separately for reference.

\begin{lem}[L3: environments and instantiation (new)]\label{lem:L3}
  If $\rho \repeq \rho'$ pointwise then
  $\evalop(t, \rho) \repeq \evalop(t, \rho')$; consequently if
  moreover $c \repeq c'$ and $w \repeq w'$ then
  $\capp_d(c, w) \repeq \capp_d(c', w')$.
\end{lem}
\begin{proof}
  For $\mathsf{mk}$-closures, $\capp$ is an $\evalop$ call on the
  common body in $\repeq$-related environments; for each closure
  combinator, $\capp$ is a fixed composite of semantic operations
  applied to the stored fields, and the fields are $\repeq$ by
  assumption.  The $\evalop$ statement is by induction on $t$: value
  formers yield equal heads with $\repeq$ arguments (interval parts
  literal by Lemma~\ref{lem:L2}); closure formers store
  $\repeq$-related environments over the same body, which is exactly
  the closure clause of Definition~\ref{def:repeq}; eliminators
  reduce to L4--L8.
\end{proof}

\begin{lem}[L4: application, projections, path application
  (new)]\label{lem:L4}
  $\repeq$ is preserved by $\mathsf{app}$, $\fst$, $\snd$,
  $\mathsf{papp}$.
\end{lem}
\begin{proof}
  Case on the common head.  $\beta$-heads ($\mathsf{lam}$,
  $\mathsf{pair}$, $\mathsf{plam}$): the operation is an
  instantiation or a field projection --- L3, or immediate.  Neutral
  heads: the spine grows by $\repeq$-related arguments, which is the
  neutral clause.  Path application at an endpoint returns a stored
  endpoint field --- immediate.
\end{proof}

\begin{lem}[L5: face forcing (new)]\label{lem:L5}
  $\repeq$ is preserved by $\mathsf{force}$ (deciding faces of
  $\hcomp$/$\glue$ values under a cofibration valuation).
\end{lem}
\begin{proof}
  The decision inspects cofibration arguments, which are literally
  equal; both sides therefore decide identically, and the surviving
  branches are $\repeq$ by assumption.
\end{proof}

\begin{lem}[L6: transport (new)]\label{lem:L6}
  If $c \repeq c'$ (family closures) and $u \repeq u'$ then
  $\transp$-evaluation on $(c, u)$ and $(c', u')$ yields
  $\repeq$-related results.
\end{lem}
\begin{proof}
  The instantiated family values $F := \capp_{d+1}(c, \ell)$ and
  $F' := \capp_{d+1}(c', \ell)$ are $\repeq$ (L3; same $\ell$ by
  L2).  The check is $\constspec$, invariant under $\repeq$ by
  Lemma~\ref{lem:readback}: \emph{both sides take the same branch}.
  Positive branch: $u \repeq u'$.  Negative branch: same head of $F,
  F'$; each structural dispatch is a fixed composite of L3--L5, L7,
  L8 and recursive L6-instances on stored fields, closing the
  bisimulation.
\end{proof}

\begin{lem}[L7: homogeneous composition (new)]\label{lem:L7}
  $\repeq$ is preserved by $\hcomp$ and by system formation.
\end{lem}
\begin{proof}
  System entries are pairs of literally-equal cofibrations and
  $\repeq$-related closures (L3); if a face is decided, both sides
  decide it identically (L5) and continue in the tube; otherwise both
  form an $\hcomp$ value with componentwise-$\repeq$ arguments.
\end{proof}

\begin{lem}[L8: $\Glue$ operations (new)]\label{lem:L8}
  $\repeq$ is preserved by $\glue$, $\unglue$ and
  $\mathsf{transpGlue}$.
\end{lem}
\begin{proof}
  $\glue$/$\unglue$: head formation and field extraction over
  componentwise-$\repeq$ data, with face decisions by L5.
  $\mathsf{transpGlue}$ is a fixed composite of instantiations,
  transports, compositions, projections and (un)gluings
  (\cite[\S6.2]{CCHM}); each step is L3--L7 or a recursive
  L6-instance.
\end{proof}

\noindent (L9 is Lemma~\ref{lem:readback} in the main text.)

\subsection{Value-level lemmas of the base component}
\label{app:valuelemmas}

By Remark~\ref{rem:samedomain} the following are statements about the
shared domain and transfer verbatim from the base theory; we record
the statements and the proof shapes for self-containment.

\begin{lem}[L10: PER laws and type invariance (B)]\label{lem:L10}
  Each $\valeq_T$ is symmetric and transitive, and if
  $T \valeq_{\U} T'$ then $\valeq_T = \valeq_{T'}$.
\end{lem}
\begin{proof}[Proof shape]
  Induction on the stratified type value.  Symmetry and transitivity
  are inherited clause-wise (e.g.\ at $\Pi$ from the codomain PERs at
  related arguments, using symmetry to swap and transitivity to
  compose through a common argument).  Type invariance: related type
  values have, clause by clause, cofinally the same decompositions
  (related domains and codomain families), so the defined PERs
  coincide; at $\U$ this is the universe induction.
\end{proof}

\begin{lem}[L11: Kripke monotonicity (B)]\label{lem:L11}
  $v \valeq_T v'$ at depth $d$ and $d' \ge d$ imply
  $v \valeq_T v'$ at depth $d'$.
\end{lem}
\begin{proof}[Proof shape]
  The quantifiers in Definition~\ref{def:valeq} range over all future
  depths; restriction is monotone.
\end{proof}

\begin{lem}[L12: conversion correctness (B)]\label{lem:L12}
  $\convop_d(v, v', T)$ accepts iff $v \valeq_T v'$.  The extension
  adds \emph{no} clause to $\convop$: conversion consumes values,
  and rule (R) only changes which values the evaluator produces.
\end{lem}
\begin{proof}[Proof shape]
  Soundness: induction on the accepting run; each clause of
  $\convop$ ($\eta$-directed at $\Pi$/$\Sigma$/$\Path$, DNF equality
  for interval arguments of neutral path applications, spine
  comparison for neutrals, componentwise elsewhere) verifies a
  defining clause of $\valeq$.  Completeness: by L13, related values
  have equal read-back, and $\convop$ accepts values with equal
  read-back (induction on the read-back).
\end{proof}

\begin{lem}[L13: read-back correctness (B)]\label{lem:L13}
  $v \valeq_T \evalop(\quo_d\,v)$; distinct normal forms of the same
  type are inconvertible; read-back is injective on
  $\valeq$-classes.
\end{lem}

\begin{lem}[L14: read-back/action commutation (new)]\label{lem:L14}
  $\quo_d(v\langle\xi\rangle) = \quo_d(v)\langle\xi\rangle$ for any
  dimension substitution $\xi$ whose domain is disjoint from the
  generic levels used by the read-back.
\end{lem}
\begin{proof}
  Induction on the read-back.  Head cases: the action preserves the
  head and acts on interval parts; read-back copies interval parts in
  DNF, and the action followed by DNF-renormalization equals
  DNF-renormalization followed by the syntactic substitution, because
  substitution into a De~Morgan expression is a homomorphism and the
  antichain DNF is canonical.  Binder cases: read-back opens the
  closure at a generic level $\notin \mathrm{dom}\,\xi$; the action
  passes into the closure environment (its definition) and the
  induction hypothesis applies to the opened body.  Neutral cases:
  spine-wise.
\end{proof}

\begin{lem}[L15: numeral inversion (B)]\label{lem:L15}
  A closed value $v$ with $v \valeq_{\mathbb{N}} v$ is a numeral.
\end{lem}
\begin{proof}[Proof shape]
  The $\valeq_{\mathbb{N}}$ clause allows numerals or neutrals with
  equal read-back; closed neutrals do not exist (a neutral's head is
  a variable level, and depth $0$ admits none).
\end{proof}

\section{The fundamental theorem, case by case}
\label{app:fund}

Throughout, $\rho \valeq_\Gamma \rho'$ abbreviates ``related
$\Gamma$-environments'', and \emph{IH} the induction hypothesis of
the outer induction (on the typing derivation).  ``Defined'' includes
termination (Definition~\ref{def:semtyping}).  Cases marked (B) are
the base component in environment form; we give their proofs in the
uniform one-paragraph style to make the claimed transposition
inspectable, with the understanding that they follow
\cite{HuberNorm} (De~Morgan) --- the correspondence for the Cartesian
layer is Appendix~\ref{app:bridge}.

\subparagraph*{Variables (B).}  $\evalop(x, \rho) = \rho(x)$, defined
and related by assumption on $\rho \valeq \rho'$.

\subparagraph*{$\Pi$-introduction (B).}  $\evalop(\lambda x.t, \rho)$
is a closure value; for the $\Pi$-clause of $\valeq$, take
$d' \ge d$ and $w \valeq_D w'$; instantiation continues as
$\evalop(t, \rho{+}[w])$, and $\rho{+}[w] \valeq \rho'{+}[w']$, so IH
for $t$ applies; definedness of the body's evaluation under every
related extension is exactly IH's definedness clause.

\subparagraph*{$\Pi$-elimination (B).}  $\evalop(t\,u, \rho) =
\mathsf{app}(\evalop(t,\rho), \evalop(u,\rho))$; by IH both parts are
defined and related; the $\Pi$-clause of $\valeq$ applied at
$w := \evalop(u,\rho) \valeq w' := \evalop(u,\rho')$ gives
definedness and relatedness of the application; the type is the
instantiated codomain, matching the typing rule by L3.

\subparagraph*{$\Sigma$ (B).}  Introduction: componentwise by IH.
Eliminations $\fst$/$\snd$: the $\Sigma$-clause is componentwise;
projections of related values are related at domain and instantiated
codomain respectively.

\subparagraph*{$\Path$-introduction (B).}
$\evalop(\langle i \rangle t, \rho)$ is a path closure; the
$\mathsf{PathP}$-clause quantifies over interval values $\iota$;
instantiation continues as $\evalop(t, \rho{+}[\iota])$ with IH; the
endpoint conditions are the typing rule's boundary premises evaluated
by IH, using L12 to convert the boundary equalities.

\subparagraph*{$\Path$-elimination (B).}  Application of related path
values to a common interval value, by the $\mathsf{PathP}$-clause;
at endpoints the stored endpoint values are returned, related by the
formation premises.

\subparagraph*{Universe rules (B).}  Codes evaluate to type values;
the $\U$-clause requires the decoded PERs to be related, which is IH
at the smaller universe level (the stratification of
Definition~\ref{def:valeq}); cumulativity is monotone inclusion of
the stratification.

\subparagraph*{$\mathbb{N}$ (B).}  $\mathsf{zero}$, $\mathsf{suc}$:
numeral clauses.  Recursor: by the $\valeq_{\mathbb{N}}$ clause the
scrutinee is a numeral or neutral with equal read-back; in the
numeral case compute by IH on the branch, by induction on the
numeral; in the neutral case both sides form neutral recursor spines
with related components, related by the neutral clause.

\subparagraph*{$\hcomp$, all formers (B).}  If a face of the system
is decided, both evaluations continue in that tube at the top of the
fill dimension (L5), related by IH on the system entry.  Otherwise:
at $\Pi$/$\Sigma$/$\Path$ the composition commutes into the
components (the base structural $\hcomp$ rules), reducing to IH and
the componentwise clauses; at base types and HITs the value is an
$\hcomp$-constructor covered by the $\hcomp$-closure clause of
$\valeq$; at $\Glue$ the composite operation is a fixed composite of
covered operations (L8).  \emph{No clause of this case mentions
transport redexes} --- the observation that lets the extension leave
it untouched.

\subparagraph*{$\Glue$ (B).}  Formation: componentwise (type values
with related base and branch data on faces).  Introduction
($\glue$): the $\Glue$-clause asks relatedness of the unglued part
and of the branch parts on their faces --- both by IH.  Elimination
($\unglue$): extraction, plus the decided-face case by L5.

\subparagraph*{HIT constructors and eliminators (B).}  For the
signatures of Section~\ref{sec:hits}: constructors are covered by
the constructor clause of $\valeq$ (value arguments by IH, interval
arguments literal); eliminators follow the recursor pattern: on
constructor values compute by IH and the signature's computation
rule; on $\hcomp$ values use the $\hcomp$-closure clause and
commute; on neutrals form related neutral spines.

\subparagraph*{Structural transport dispatches (B), and where the
inner induction runs.}  These are the negative-branch continuations
of the transport clause; we spell out the two used in
Section~\ref{sec:switchover} and the pattern for the rest.
\begin{itemize}
\item \emph{$\Pi$-line.}  The result is the closure value
  $\lambda x_1.\transp(\mathrm{cod\ line})\,
  (f\,(\transp(\mathrm{dom\ line})^{-}\,x_1))$.  To place it in the
  $\Pi$-clause at the endpoint type, take related arguments
  $w \valeq w'$ at the endpoint domain; the reversed domain transport
  is a transport along a line whose fibers are \emph{structurally
  smaller type values} than the $\Pi$-line, so the \emph{inner}
  induction (on type values, by which $\valeq$ is defined) provides
  its definedness and relatedness; likewise the codomain transport.
  The outer IH is used only for the premises ($f$, the line).  This
  is the division of labor claimed in
  Theorem~\ref{thm:fundamental}: the evaluator's recursion along the
  line's structure is mirrored by the inner induction, never by the
  outer one.
\item \emph{$\Sigma$-line, $\Path$-line, HIT lines, $\Glue$ line.}
  Same pattern: each inner transport or composition is at a
  structurally smaller type value (domain fiber, codomain fiber over
  a transported point, tube line, branch line), discharged by the
  inner induction; $\hcomp$ components by the $\hcomp$ case;
  $\mathsf{transpGlue}$'s steps by L8 and the inner induction.
\item \emph{Neutral line.}  Both sides form related neutral
  transport spines (neutral clause).
\end{itemize}

\subparagraph*{Transport, positive branch (new).}  Suppose
$\constspec(F, \ell)$.  By Theorem~\ref{thm:checksound} the endpoint
type values have equal read-back and decode to the same PER; by IH
$\evalop(u,\rho) \valeq \evalop(u,\rho')$ at the source endpoint,
hence at the target endpoint \emph{by identity of the PERs} --- no
closure property under judgmental equality is invoked.  Definedness
is IH's.  By Lemma~\ref{lem:subst} and Lemma~\ref{lem:readback} the
branch taken is a function of the $\repeq$-class, so the case is
stable under the environment/substitution manipulations used
elsewhere in the proof.

\subparagraph*{Congruence lemmas (Theorem~\ref{thm:admissible}(2)).}
For each former, the lemma ``subterms related $\implies$ whole
related'' is the corresponding case above read with the IH replaced
by the assumption; we list the non-obvious instances:
$\transp$-congruence needs the branch-agreement argument (the two
families are $\valeq$, hence --- being values of well-typed families
--- of equal read-back by L13, so $\constspec$ agrees);
$\hcomp$-congruence needs L5 to align face decisions; HIT-eliminator
congruence needs the neutral clause when the scrutinee is neutral.
All others are componentwise.

\section{Checker exactness: the full clause table}
\label{app:exact}

Theorem~\ref{thm:exact} rests on the claim that \texttt{usesLvl} and
$\quo$ have mirrored recursion structures.  This appendix makes the
claim inspectable.  Throughout, ``$\ell \in \mathsf{iv}(r)$'' means
the level $\ell$ occurs in the interval value $r$, and we use that
interval quotation is \emph{structural} (it maps the De~Morgan
operators of the value to the corresponding term formers), so
$\ell \in \mathsf{iv}(r) \iff \ell \in \FV(\quo(r))$.

\subparagraph*{Definedness transfer.}  Both functions begin by
forcing the value's head under the current cofibration valuation ---
the same $\mathsf{force}$ --- and recurse along the same
sub-computations: on value arguments structurally, and on closures
via the \emph{same} instantiation $\capp$ at the \emph{same} fresh
level (equal to the current depth, Lemma~L2).  Hence the big-step
derivation trees of $\texttt{usesLvl}(\ell, F)$ and $\quo(F)$ are
isomorphic --- one is defined iff the other is --- with two
exceptions in \texttt{usesLvl}'s favour, both of which only prune
sub-derivations: the level-bound cut and the early exit of the
disjunctive traversal (\texttt{usesLvl} may stop at the first
occurrence).  Pruning cannot destroy definedness.

\subparagraph*{Clause table.}  Value cases (heads after
$\mathsf{force}$):

\begin{center}
\footnotesize
\begin{tabular}{lll}
  head & \texttt{usesLvl} & $\quo$ \\\hline
  constants ($\U_k$, $\mathbb{N}$, $0$, \dots) &
    $\mathsf{false}$ & closed term \\
  $\mathsf{suc}\,n$, injections, $\mathsf{tin}$, $\mathsf{qin}$,
    \dots & recurse on arguments & recurse on arguments \\
  HIT cells (e.g.\ $\mathsf{merid}\,a\,r$,
    $\mathsf{emloop}\,g\,r$, squash cells) &
    args $\lor$ $\ell \in \mathsf{iv}(r_i)$ &
    args, $\quo(r_i)$ \\
  $\Pi$/$\Sigma$ (domain, codomain closure) &
    dom $\lor$ binder-inst.\ & dom, binder-inst.\ \\
  $\lambda$-value (closure) & binder-inst.\ & binder-inst.\ \\
  $\mathsf{PathP}$ (line closure, endpoints) &
    i-binder-inst.\ $\lor$ endpoints &
    i-binder-inst., endpoints \\
  path abstraction (closure) & i-binder-inst.\ &
    i-binder-inst.\ \\
  interval value $r$ & $\ell \in \mathsf{iv}(r)$ & $\quo(r)$ \\
  $\Glue$ type / $\glue$ value &
    components $\lor$ cof mentions & components, cof terms \\
  neutral & spine (below) & spine (below) \\
\end{tabular}
\end{center}

Neutral cases: variable ($\ell$-equality of the head level
$\leftrightarrow$ occurrence of the corresponding de~Bruijn index in
the quoted variable); application, projections, path application,
recursors (spine pointwise, motive closures by binder
instantiation); \emph{neutral transport} (family by interval-binder
instantiation, argument structurally); \emph{neutral $\hcomp$} (type
and base structurally, each tube by cofibration mentions plus
interval-binder instantiation of the tube closure --- the case that
the legacy dead code of Remark~\ref{rem:archaeology} would have
treated structurally); $\unglue$ and the HIT eliminators (spine
pointwise).  In every row the right column's free-variable content
is computed by the same recursion as the left column's disjunction,
so the biconditional propagates by induction on the (isomorphic)
derivations.

\subparagraph*{Completeness of the level-bound cut.}  The invariant
of Lemma~L1 is maintained by \emph{every} smart constructor,
including the closure formers: the bound of a closure is the maximum
of the bounds of its stored environment entries and combinator
fields.  If $\ell \ge \mathrm{lb}(v)$, then $\ell$ is not among
$v$'s stored levels; levels introduced by instantiations during
read-back are fresh, i.e.\ $\ge$ the current depth $> \ell$'s
binding depth, and are bound by the binders the read-back produces.
So $\ell \notin \FV(\quo(v))$: the cut answers exactly what the
specification would.

\subparagraph*{Transparency of memoization.}  The memo hooks cache
the result of $\texttt{usesLvl}(\ell, \cdot)$ keyed by the level and
the \emph{physical identity} of the value (pointer equality).
Physical identity implies structural identity, and
$\texttt{usesLvl}$ is a pure function of its arguments (it neither
allocates observable state nor depends on any); caching a pure
function at a sound key is observationally transparent.  The
generation-stamp discipline ensures a value reachable from two
environments is the \emph{same} value, so pointer keys do not
conflate distinct values.  Likewise the depth argument is determined
by the value's stamp, so omitting it from the key loses nothing.

\section{Discharging the base component for the minimal core}
\label{app:core}

This appendix proves component (B) --- and hence
Theorems~\ref{thm:fundamental}, \ref{thm:admissible}
and~\ref{thm:canonicity} unconditionally
(Corollary~\ref{thm:core}) --- for the \emph{minimal core}
\[
  T^{\mathrm{core}} \;:=\;
  \Pi,\ \Sigma,\ \mathsf{PathP},\ \mathbb{N}
  \ \text{with}\ \transp\ \text{and}\ \hcomp
  \quad(\text{no } \Glue,\ \text{no universes},\ \text{no HITs}),
\]
over the De~Morgan interval; the Cartesian variant is a
simplification of the same proof (drop the connection cases of the
interval algebra).  The fragment is chosen as the \emph{smallest}
one in which $\J$ is definable (Proposition~\ref{prop:J}); the
extension by $S^1$ is treated separately, as a proposition with a
proof sketch and artifact evidence
(Section~\ref{app:s1ext}), precisely so that the headline theorem
is not entangled with higher-inductive proof obligations.

\subsection{Types, values, normal forms, and the induction order}
\label{app:coretypes}

Since there is no universe, types are a separate syntactic
category:
\[
  T, S ::= \mathbb{N} \mid \Pi x{:}T.S \mid \Sigma x{:}T.S \mid
  \mathsf{PathP}\,(\langle i \rangle T)\,t\,u .
\]

\begin{lem}[Type-value inversion]\label{lem:coreinv}
  If $\evalop(T, \rho)$ is defined, its head is the value former
  matching $T$'s head former, and its closure fields store the
  syntactic subexpressions of $T$ with (extensions of) $\rho$:
  $\evalop(\Pi x{:}T.S, \rho) = \mathsf{v}\Pi(\evalop(T,\rho),
  \langle \rho; S \rangle)$, and similarly for $\Sigma$ and
  $\mathsf{PathP}$.
\end{lem}
\begin{proof}
  Immediate from the evaluator's type clauses: each is a single
  head-formation step storing a $\mathsf{mk}$-closure.
\end{proof}

\begin{rem}[The induction order, precisely]\label{rem:order}
  Lemma~\ref{lem:coreinv} is what makes the fragment's metatheory
  elementary: every instantiation of a stored closure of a type
  value is an evaluation of a \emph{proper syntactic subexpression}
  in an extended environment.  Accordingly, the logical relation
  below is defined by \emph{structural recursion on the syntactic
  type $T$}, carrying the environment as a parameter --- when the
  $\Pi$-clause speaks of $\capp(C, w)$, it is, by
  Lemma~\ref{lem:coreinv}, speaking of $\evalop(S, \rho{+}[w])$
  with $S$ a proper subterm of $\Pi x{:}T.S$.  This is the exact
  well-founded order behind every ``structurally smaller type
  value'' step in this appendix; no ordinal, size measure on
  values, or universe stratification appears anywhere.
\end{rem}

\noindent Normal and neutral forms of the core (the target of
read-back):
\begin{align*}
  \mathsf{nf} &::= \mathsf{ze} \mid \mathsf{su}\,\mathsf{nf}
    \mid \lambda x.\mathsf{nf} \mid (\mathsf{nf}, \mathsf{nf})
    \mid \langle i \rangle \mathsf{nf} \mid \mathsf{ne} \\
  \mathsf{ne} &::= x \mid \mathsf{ne}\,\mathsf{nf}
    \mid \fst \mathsf{ne} \mid \snd \mathsf{ne}
    \mid \mathsf{ne}\,r
    \mid \mathsf{natrec}(\mathsf{nf}, \mathsf{nf}, \mathsf{nf},
      \mathsf{ne})
    \mid \transp^i \mathsf{T}_{\mathsf{ne}}\,\varphi\,\mathsf{nf}
    \mid \hcomp\,[\vec\varphi \mapsto \vec{\mathsf{nf}}]\,
      \mathsf{nf}
\end{align*}
where in the core the transport-neutral case does not arise (all
core type values are canonical, Lemma~\ref{lem:coreinv}), path
application $\mathsf{ne}\,r$ is restricted to $r$ not an endpoint
(the $\eta$/boundary rules fire otherwise), and $\hcomp$-neutrals
occur only at $\mathbb{N}$ over a neutral base with an undecided
system (Lemma~\ref{lem:corehcompN}).  Read-back targets exactly
this grammar, clause by clause of $\quo$
(Appendix~\ref{app:exact}).

\subsection{Worlds: depths, dimensions, and faces}
\label{app:coreworlds}

Contexts of the calculus contain term variables, \emph{interval
variables} and \emph{cofibration restrictions} $\Gamma, \varphi$;
and the boundary conditions of $\hcomp$ hold only \emph{under
faces}.  The logical relation must therefore be indexed not by
depths alone but by:

\begin{defi}[Worlds]\label{def:worlds}
  Worlds are defined relative to a fixed \emph{sorted context
  shape}: the context determines, for each level $< d$, whether it
  is a term level or an interval level, and all notions below
  respect this sorting (the implementation records it in the
  context; values never confuse the sorts, by the typing of the
  smart constructors).  A \emph{world} is then a pair
  $w = (d, \psi)$ of a depth $d$ (making available the term and
  interval levels $< d$, as sorted by the context shape) and a
  cofibration $\psi$ over the interval levels $< d$ (the conjunction of the
  faces currently assumed; $\psi = \ii$ is the empty assumption,
  $\psi = \io$ the absurd one).  A \emph{world map}
  $(d', \psi') \to (d, \psi)$ is a dimension substitution
  $\xi$ from the interval levels $< d$ to $\DM{<d'}$ (extended by
  the identity on term levels, with $d \le d'$ on the term part)
  such that $\psi' \models \psi\langle\xi\rangle$.  Weakenings
  ($d \le d'$, $\xi = \mathrm{id}$, face strengthening
  $\psi' \models \psi$) and face restrictions
  ($\psi' := \psi \wedge \varphi$) are the two generating cases
  used below.
\end{defi}

\begin{defi}[Semantic context satisfaction]\label{def:semctx}
  Related environments at a context, at a world $w = (d, \psi)$:
  term variables pointwise related (at the types of the context,
  each in the preceding environments, at $w$); interval variables
  assigned \emph{the same} interval value over the levels $< d$ on
  both sides (Definition~\ref{def:repeq}'s literal-interval
  discipline); and for a restriction entry $\varphi$ of the
  context, $\psi \models \varphi[\rho]$ --- the world must satisfy
  the restriction.
\end{defi}

At the absurd world ($\psi = \io$, i.e.\ $\psi \models \bot$) all
values of all types are related by convention; this is what makes
empty-intersection compatibility conditions (such as
\eqref{eq:P3}) genuinely vacuous rather than informally so.

\subsection{The logical relation}
\label{app:corerel}

For each syntactic type $T$, world $w = (d, \psi)$, and pair of
environments $\rho, \rho'$ related at $T$'s context at $w$
(Definition~\ref{def:semctx}), a relation
$v \sim^w_{T[\rho,\rho']} v'$ on values, by structural recursion on
$T$; every clause is implicitly \emph{Kripke over worlds}: it
quantifies over all $w' \to w$, with the values transported by the
corresponding dimension action, and all clauses are Kleene (the
mentioned applications and instantiations are required to be
defined):

\begin{itemize}
\item $\mathbb{N}$: both force to the same numeral, or to neutrals
  with equal read-back (including stuck $\hcomp$-neutrals over
  neutral bases);
\item $\Pi x{:}T.S$: for all worlds $w' \to w$ and all
  $u \sim^{w'}_{T[\rho,\rho']} u'$:
  $f\,u \sim^{w'}_{S[\rho+[u],\, \rho'+[u']]} f'\,u'$;
\item $\Sigma x{:}T.S$: $\fst u \sim^w_{T[\rho,\rho']} \fst u'$
  and
  $\snd u \sim^w_{S[\rho+[\fst u],\, \rho'+[\fst u']]} \snd u'$;
\item $\mathsf{PathP}\,(\langle i \rangle T)\,l\,r$: for every
  world $w' \to w$ and interval value $\iota$ over $w'$:
  $p\,\iota \sim^{w'}_{T[\rho+[\iota],\, \rho'+[\iota]]} p'\,\iota$;
  at $\iota = \io$ (resp.\ $\ii$) the boundary rules force the
  values of $l$ (resp.\ $r$), so the endpoint conditions are
  instances of the same clause.
\end{itemize}
A \emph{computable type in $\rho$ at $w$} is one all of whose
stored components are so related (for $\mathsf{PathP}$: the line,
at every instantiation over every future world, and both
endpoints).

\begin{defi}[Related systems]\label{def:relsys}
  Systems $[\varphi_k \mapsto t_k]$, $[\varphi_k \mapsto t_k']$
  (same cofibrations, by Definition~\ref{def:repeq}) are related at
  $T$, $w = (d, \psi)$ if: for each $k$, the tubes are related
  \emph{at the restricted world} $(d, \psi \wedge \varphi_k)$, at
  every instantiation of their interval binder; and for each pair
  $k, l$, the tubes agree (are related to each other) at
  $(d, \psi \wedge \varphi_k \wedge \varphi_l)$.  A base $u_0$ is
  \emph{adapted} to the system if
  $u_0 \sim^{(d, \psi \wedge \varphi_k)} t_k(\io)$ for each $k$ ---
  the semantic boundary condition.
\end{defi}

\begin{lem}[PER laws, environment irrelevance, world
  functoriality]\label{lem:coreper}
  Each $\sim^w_T$ is symmetric and transitive; replacing $\rho,
  \rho'$ by environments related to them yields the same relation;
  and relatedness is \emph{functorial in the world}: for every
  world map $\theta : w' \to w$ with dimension action
  $\langle\theta\rangle$,
  $v \sim^w_T v' \implies
  v\langle\theta\rangle \sim^{w'}_{T} v'\langle\theta\rangle$
  (\emph{face restriction / weakening lemma}).
\end{lem}
\begin{proof}
  Structural induction on $T$.  PER laws: $\mathbb{N}$ by equality
  of numerals and read-backs; $\Pi$ by swapping and composing
  through a common argument; $\Sigma$ componentwise with
  environment irrelevance at $S$; $\mathsf{PathP}$ pointwise.
  Environment irrelevance is inherited from the clauses'
  quantification.  Functoriality: every clause quantifies over all
  future worlds, and world maps compose, so the quantified
  conditions restrict; the only content is at the base cases ---
  the dimension action commutes with forcing (a face decided at
  $w$ stays decided at $w'$ since
  $\psi' \models \psi\langle\theta\rangle$; new decisions can only
  collapse both sides identically, L5), and read-back commutes with
  the action (Lemma~L14) --- so numeral/neutral disjuncts are
  preserved.  At the absurd world the convention closes every case.
\end{proof}

\begin{lem}[Boundary preservation]\label{lem:boundary}
  Let the system and adapted base be related at $T$, $w = (d,
  \psi)$ (Definition~\ref{def:relsys}), and suppose the $\hcomp$
  at $T$ is defined (as it is in each case of
  Lemmas~\ref{lem:corehcompN}--\ref{lem:corehcompstr}).  Then at
  each restricted world $(d, \psi \wedge \varphi_k)$, the composite
  is related to the tube top $t_k(\ii)$.
\end{lem}
\begin{proof}
  At $(d, \psi \wedge \varphi_k)$ the face $\varphi_k$ is
  satisfied, so forcing decides it and the composite collapses to
  $t_k(\ii)$ on both sides (the $\hcomp$ face rule; L5); the tube
  tops are related there by Definition~\ref{def:relsys}.
  Functoriality (Lemma~\ref{lem:coreper}) transports the ambient
  relatedness data to the restricted world, so the collapse happens
  in a coherent context.
\end{proof}

\subsection{Reflect and reify, former by former}
\label{app:corereify}

Reflect and reify are proved by \emph{one} mutual structural
induction on the syntactic type: the lemma for a former may use both
lemmas at that former's proper subexpressions.  We state them per
former, as separate lemmas, so that each obligation is inspectable
in isolation; the induction discipline is: L(T) may invoke
L$'$($T'$) for any $T'$ a proper subterm of $T$ and any L$'$ among
the eight lemmas below.

\begin{lem}[Reflect at $\mathbb{N}$]\label{lem:reflN}
  Neutrals $n, n'$ of type $\mathbb{N}$ with $\quo(n) = \quo(n')$
  satisfy $n \sim_{\mathbb{N}} n'$.
\end{lem}
\begin{proof}
  The neutral disjunct of the $\mathbb{N}$-clause, verbatim.
\end{proof}

\begin{lem}[Reify at $\mathbb{N}$]\label{lem:reifN}
  $v \sim_{\mathbb{N}} v' \implies \quo(v) = \quo(v')$, both
  defined.
\end{lem}
\begin{proof}
  Same-numeral case: numerals quote to themselves.  Neutral case:
  the clause \emph{is} read-back equality.
\end{proof}

\begin{lem}[Reflect at $\Pi$]\label{lem:reflPi}
  Let $\Pi x{:}T.S$ be computable in $\rho, \rho'$ and $n, n'$
  neutrals of that type with equal read-back.  Then
  $n \sim_{\Pi x{:}T.S\,[\rho,\rho']} n'$.
\end{lem}
\begin{proof}
  Given $d' \ge d$ and $w \sim_{T} w'$: by reify at the subterm
  $T$ (the mutual-induction instance of this lemma family),
  $\quo(w) = \quo(w')$, so the applications $n\,w$, $n'\,w'$ are
  neutrals with equal read-back
  $\quo(n)\,\quo(w) = \quo(n')\,\quo(w')$; reflect at the subterm
  $S$ (in the extended, related environments) relates them at $S$,
  which is the $\Pi$-clause's requirement.  Here and below,
  ``reflect/reify at $T$'' always abbreviates the mutual-induction
  instance at the proper subterm $T$.
\end{proof}

\begin{lem}[Reify at $\Pi$]\label{lem:reifPi}
  $f \sim_{\Pi x{:}T.S\,[\rho,\rho']} f' \implies
  \quo(f) = \quo(f')$, both defined.
\end{lem}
\begin{proof}
  Let $x_d$ be the generic neutral of type $T$ at the current
  depth; it is related to itself by reflect at $T$ (its read-back
  is the variable).  The $\Pi$-clause gives
  $f\,x_d \sim_{S[\rho+[x_d]]} f'\,x_d$, defined; reify at $S$
  gives equal body read-backs; $\quo$ at $\Pi$-type values produces
  the $\lambda$-abstraction of exactly that body
  (Appendix~\ref{app:exact}), so the read-backs agree.
\end{proof}

\begin{lem}[Reflect at $\Sigma$]\label{lem:reflSig}
  Under the same hypotheses at $\Sigma x{:}T.S$:
  $n \sim_{\Sigma x{:}T.S} n'$.
\end{lem}
\begin{proof}
  $\fst n, \fst n'$ are neutrals with equal read-back, related at
  $T$ by reflect at $T$; hence the second-component type
  $S[\rho{+}[\fst n]]$ is computable
  (Lemma~\ref{lem:coreper}, environment irrelevance), and
  $\snd n, \snd n'$ are neutrals with equal read-back, related at
  it by reflect at $S$.  Both components together are the
  $\Sigma$-clause.
\end{proof}

\begin{lem}[Reify at $\Sigma$]\label{lem:reifSig}
  $u \sim_{\Sigma x{:}T.S} u' \implies \quo(u) = \quo(u')$.
\end{lem}
\begin{proof}
  Componentwise by reify at $T$ and at $S$ (the latter in the
  environments extended by the related first components); $\quo$
  produces the pair of the component read-backs.
\end{proof}

\begin{lem}[Reflect at $\mathsf{PathP}$]\label{lem:reflPath}
  Let $P := \mathsf{PathP}(\langle i \rangle T)\,l\,r$ be
  computable in $\rho, \rho'$ and $n, n'$ neutrals of type $P$ with
  equal read-back.  Then $n \sim_P n'$.
\end{lem}
\begin{proof}
  Take any interval value $\iota$.  If $\iota$ is not forced to an
  endpoint, $n\,\iota$, $n'\,\iota$ are neutral path applications;
  their read-backs agree because interval read-back is structural
  and the spine read-backs agree; reflect at $T$ (in the
  environments extended by $\iota$) relates them.  If
  $\iota = \io$ (resp.\ $\ii$), the boundary rule returns the
  stored endpoint values of $l$ (resp.\ $r$) on both sides ---
  related because $P$ is computable (its endpoint components are
  related by definition of computable type).  Both cases are what
  the $\mathsf{PathP}$-clause demands.
\end{proof}

\begin{lem}[Reify at $\mathsf{PathP}$]\label{lem:reifPath}
  $p \sim_P p' \implies \quo(p) = \quo(p')$.
\end{lem}
\begin{proof}
  Instantiate at the generic interval level $\iota_d$: related
  bodies at $T[\rho{+}[\iota_d]]$ by the clause; reify at $T$
  gives equal body read-backs; interval arguments of any neutral
  path applications occurring inside compare equal by canonical
  antichain-DNF (sound and complete for the free De~Morgan
  algebra); $\quo$ produces the path abstraction of the common
  body.
\end{proof}

\noindent The mutual induction across the eight lemmas is
well-founded by Remark~\ref{rem:order}: every cross-reference is at
a proper syntactic subterm.

\subsection{Conversion correctness, restructured}
\label{app:coreconv}

Conversion correctness splits into a purely \emph{syntactic}
characterization and two \emph{type-level} semantic lemmas.  The
split is deliberate: read-back equality and quote--eval idempotence
are different theorems, and the load-bearing semantic content is
needed \emph{only at type values}, where the core's type-value
inversion (Lemma~\ref{lem:coreinv}) makes it provable by structural
induction on syntax --- the generic-to-arbitrary passage never has
to be performed on arbitrary function \emph{values}.

\begin{lem}[Conversion is read-back equality]\label{lem:convrb}
  For values of a core type, $\convop(v, v', T\text{-value})$
  accepts iff $\quo(v) = \quo(v')$, and the accepting run is
  defined iff both read-backs are.
\end{lem}
\begin{proof}
  Purely syntactic, no logical relation involved.
  ($\Rightarrow$) Induction on the accepting run: each $\convop$
  clause compares exactly what the corresponding $\quo$ clause
  emits --- the $\eta$-directed $\Pi$/$\Sigma$/$\mathsf{PathP}$
  clauses instantiate at the same generic levels as
  $\quo$'s binder cases; numerals, neutral spines and interval
  arguments (DNF) coincide with their quotations.
  ($\Leftarrow$) Induction on the common read-back along the
  normal-form grammar of Section~\ref{app:coretypes}: each
  production is matched by exactly one $\convop$ clause.
  Definedness transfers as in Theorem~\ref{thm:exact}: the two
  traversals visit the same instantiations.
\end{proof}

\noindent Lemma~\ref{lem:convrb} already yields, with no semantic
input, that $\aconv$ restricted to the core is an equivalence
relation and a congruence in all value-forming positions: it is
equality of read-backs.

\begin{lem}[Type-level idempotence]\label{lem:typeidem}
  Let $F = \evalop(T, \rho)$ be a computable type value at $w$.
  Then $\quo(F)$ is a syntactic core type, it is well-typed in the
  generic context of $w$, and $\evalop(\quo(F), \rho_{\mathrm{gen}})$
  is defined and determines \emph{the same relation} as $F$:
  $\sim_{\quo(F)[\rho_{\mathrm{gen}}]} \;=\; \sim_{T[\rho]}$.
\end{lem}
\begin{proof}
  Structural induction on $T$, using type-value inversion
  throughout.  $\mathbb{N}$: $\quo(F) = \mathbb{N}$, evaluation
  returns $\mathsf{v}\mathbb{N}$, same clause.
  $\Pi x{:}T.S$: $\quo(F) = \Pi x{:}\quo(D).\quo(C\,x_d)$ where $D$
  is the domain value and $C\,x_d$ the codomain instantiated at the
  generic level; by the induction hypothesis at $T$ the quoted
  domain re-evaluates to a type value with the same relation, and
  for every $u \sim u'$ at a future world the quoted codomain
  re-evaluates over $\rho_{\mathrm{gen}}{+}[u]$ --- by
  Lemma~\ref{lem:coreinv} both the original and the re-evaluated
  codomain instantiations are evaluations of the \emph{syntactic}
  $S$ (resp.\ its read-back) in related environments, so the
  induction hypothesis at $S$ applies pointwise.  $\Sigma$,
  $\mathsf{PathP}$: the same pattern (for $\mathsf{PathP}$,
  pointwise over interval values at all future worlds, endpoints by
  the induction hypothesis at the endpoint type).  Well-typedness
  of the read-back is the same induction read syntactically.
\end{proof}

\begin{lem}[Reindexing at $\mathsf{PathP}$]\label{lem:reindex}
  Let $L, L'$ be line closures of core type (over the syntactic
  types $T, T'$, by inversion) and suppose
  $\evalop(T, \rho{+}[\iota_d]) $ and
  $\evalop(T', \rho'{+}[\iota_d])$ determine the same relation at
  the generic interval level $\iota_d$ (at all future worlds).
  Then for \emph{every} interval value $\iota$ over any future
  world, $\evalop(T, \rho{+}[\iota])$ and
  $\evalop(T', \rho'{+}[\iota])$ determine the same relation.
\end{lem}
\begin{proof}
  World functoriality (Lemma~\ref{lem:coreper}): the map that sends
  the generic level to $\iota$ is a world map
  (Definition~\ref{def:worlds}), the dimension action transports
  the generic instance to the $\iota$-instance, and evaluation
  commutes with the action on the interval entry of the environment
  (the action acts pointwise on environments,
  Section~\ref{sec:rawdomain}).  Note this is a \emph{dimension}
  substitution --- available as a world map --- not a term
  substitution; no appeal to Lemma~\ref{lem:subst} is needed.
\end{proof}

\begin{lem}[Type conversion soundness]\label{lem:typeconvsound}
  Let $F, F'$ be computable type values at $w$.  If
  $\convop(F, F', \U)$ accepts, then $F$ and $F'$ determine the
  same relation.
\end{lem}
\begin{proof}
  By Lemma~\ref{lem:convrb}, $\quo(F) = \quo(F')$; by
  Lemma~\ref{lem:typeidem} both $F$ and $F'$ determine the same
  relation as the evaluation of that common syntactic type in the
  generic environment.  (Unfolded, the generic-to-arbitrary passage
  happens inside Lemma~\ref{lem:typeidem}'s $\Pi$ case by
  structural induction on the \emph{syntactic} codomain, and at
  $\mathsf{PathP}$ by Lemma~\ref{lem:reindex} --- never on
  arbitrary function values.)
\end{proof}

\begin{lem}[Conversion completeness]\label{lem:coreconvcomp}
  If $v \sim_{T[\rho,\rho']} v'$ then
  $\convop(v, v', T\text{-value})$ accepts.
\end{lem}
\begin{proof}
  By the reify lemmas
  (\ref{lem:reifN}, \ref{lem:reifPi}, \ref{lem:reifSig},
  \ref{lem:reifPath}) the read-backs are equal;
  Lemma~\ref{lem:convrb} ($\Leftarrow$) concludes.
\end{proof}

\begin{rem}[What happened to term-level semantic soundness]
  \label{rem:nosemconv}
  The development never needs ``$\convop$ accepts $\Rightarrow$
  related'' for arbitrary function \emph{values}: the fundamental
  theorem's conversion case switches \emph{type} predicates
  (Lemma~\ref{lem:typeconvsound}); the admissibility package uses
  the equivalence/congruence structure of $\aconv$, which
  Lemma~\ref{lem:convrb} provides syntactically; and the constancy
  argument uses idempotence at \emph{type} values
  (Lemma~\ref{lem:typeidem}).  We therefore do not claim the
  term-level statement, and no proof in this paper depends on it.
  This is the precise resolution of the apparent
  generic-to-arbitrary gap: at types it is closed by structural
  induction on syntax; at terms it is never used.
\end{rem}

\subsection{$\hcomp$ computability, former by former}
\label{app:corehcomp}

\noindent\emph{Induction discipline for the Kan lemmas.}
Lemmas~\ref{lem:corehcompN}--\ref{lem:coretransp} are proved by
\emph{one} simultaneous structural induction on the syntactic type:
the $\Sigma$-case of composition invokes transport at the subterm
$S$, and the $\mathsf{PathP}$-case of transport invokes composition
at the subterm $T$ --- every cross-reference between the two lemmas
is at a proper syntactic subterm, so the joint induction is
well-founded by Remark~\ref{rem:order}.  We keep the statements
separate for readability.

\begin{lem}[$\hcomp$ at $\mathbb{N}$]\label{lem:corehcompN}
  If the system is related and the base adapted, at $w$
  (Definition~\ref{def:relsys}), the $\hcomp$s at $\mathbb{N}$ are
  defined and related at $w$, with the boundary behaviour of
  Lemma~\ref{lem:boundary}.
\end{lem}
\begin{proof}
  If some face is decided true, both sides return that tube's top
  (L5), related by assumption.  Otherwise, by cases on the (common,
  by the $\mathbb{N}$-clause) shape of the base.  \emph{Numeral
  base:} by induction on the numeral, using the kernel's
  computation rules --- an $\hcomp$ at $\mathbb{N}$ with base
  $\mathsf{ze}$ (resp.\ $\mathsf{su}\,m$) peels the tubes through
  the constructor ($\mathsf{ze}$/$\mathsf{su}$ commute with
  $\hcomp$): tube tops over a numeral base at undecided faces are
  themselves related numeral-shaped values by the boundary
  condition, and the peeled composite recurses at $m$.
  \emph{Neutral base:} both sides form the stuck
  $\hcomp$-\emph{neutral} with equal read-back (tubes reified by
  the reify lemmas under their interval binder; equal
  cofibrations), covered by the neutral part of the
  $\mathbb{N}$-clause.  \emph{Closed instance remark} (used for
  canonicity): a closed cofibration contains no interval
  variables, hence is decided; so closed $\hcomp$s at $\mathbb{N}$
  always take the decided-face or numeral path and compute to
  numerals --- no closed stuck composite exists.
\end{proof}

\begin{lem}[$\hcomp$ at $\Pi$, $\Sigma$, $\mathsf{PathP}$]
  \label{lem:corehcompstr}
  Same hypotheses (Definition~\ref{def:relsys}, at $w$); the
  structural composites are defined and related at $w$, with the
  boundary behaviour of Lemma~\ref{lem:boundary}.
\end{lem}
\begin{proof}
  \emph{$\Pi$:} the composite is the closure value
  $\lambda w.\,\hcomp\,[\vec\varphi \mapsto \vec t\,w]\,(u_0\,w)$;
  applying to related $w \sim_T w'$ gives related tubes and base at
  $S[\rho{+}[w]]$ (the $\Pi$-clause, pointwise), and the induction
  hypothesis of the outer structural recursion on $T$ (the
  composite at $S$) closes the case.
  \emph{$\Sigma$:} let the input be
  $\hcomp^{\,\iota}\,(\Sigma x{:}T.S)\,
  [\vec\varphi_k \mapsto t_k]\,u_0$ in environment $\rho$.  The
  first components compose at $T$:
  \[
    c \;:=\; \hcomp^{\,\iota}\,T\,
      [\vec\varphi_k \mapsto \fst t_k]\,(\fst u_0)
    \;:\; T[\rho],
  \]
  and the kernel forms their \emph{filler}, the truncation of the
  same composition at $\iota$:
  \[
    \overline{c}(\iota) \;:=\; \hcomp^{\,\iota'}\,T\,
      [\vec\varphi_k \mapsto \fst t_k(\iota \wedge \iota'),\
       (\iota{=}0) \mapsto \fst u_0]\,(\fst u_0)
    \;:\; T[\rho],
  \]
  with $\overline{c}(0) = \fst u_0$ and $\overline{c}(1) = c$,
  related to its primed counterpart at every $\iota$ by this lemma
  at $T$ (the truncated system is again a related system).  The
  second components then compose \emph{heterogeneously along the
  line $\iota \mapsto S[\rho{+}[\overline{c}(\iota)]]$}:
  \[
    \snd(\text{result}) \;:=\;
    \vcomp^{\,\iota}\,\bigl(S[\overline{c}(\iota)/x]\bigr)\,
      [\vec\varphi_k \mapsto \snd t_k]\,(\snd u_0)
    \;:\; S[\rho{+}[c]],
  \]
  whose transport part is along instantiations of the proper
  subexpression $S$ at the related environments
  $\rho{+}[\overline{c}(\iota)]$ (Lemma~\ref{lem:coretransp} at
  $S$), and whose $\hcomp$ part is the induction hypothesis at $S$.
  The boundary premises are inherited: on $\varphi_k$, the pair
  $t_k$ has $\fst t_k$ matching the filler on that face
  ($\overline{c}(\iota) = \fst t_k(\iota)$ there, by the filler's
  face rule), so $\snd t_k$ is a legitimate tube for the
  heterogeneous composite; at $\iota = 0$ the composite starts at
  $\snd u_0 : S[\rho{+}[\fst u_0]] =
  S[\rho{+}[\overline{c}(0)]]$.  The resulting pair is related at
  $\Sigma x{:}T.S$ by the $\Sigma$-clause, whose second-component
  environment is $\rho{+}[c]$, exactly the filler at $1$.  All
  instantiations are of proper subexpressions
  (Remark~\ref{rem:order}).
  \emph{$\mathsf{PathP}$:} the composite at
  $\mathsf{PathP}(\langle i \rangle T)\,l\,r$ is
  $\langle j \rangle\, \hcomp\,[\vec\varphi \mapsto \vec t\,j,\
  j{=}0 \mapsto l,\ j{=}1 \mapsto r]\,(u_0\,j)$: the system is
  \emph{extended with the endpoint tubes}.  Endpoint coherence ---
  the reviewer-facing detail --- is exactly the boundary check: at
  $j = \io$ the added face is decided true and the composite
  collapses to $l$ (the $\hcomp$ face rule), matching the
  $\mathsf{PathP}$-clause's endpoint condition; likewise at $\ii$;
  at generic $j$ the composite is the induction hypothesis at $T$
  with the enlarged (still pairwise-compatible: the original tubes
  agree with $l, r$ at the endpoints by \emph{their} boundary
  conditions) system.
\end{proof}

\subsection{Prioritized transport computability, former by former}
\label{app:coretransp}

\begin{lem}\label{lem:coretransp}
  If the family closure (over a syntactic core type in one
  interval binder) and the argument are related, the prioritized
  transport is defined and related at the endpoint type.
\end{lem}
\begin{proof}
  \emph{Positive branch:} Theorem~\ref{thm:checksound}; its
  ingredients --- read-back/action commutation (Lemma~L14) and
  type-level idempotence --- are available for the core as
  Lemma~\ref{lem:typeidem}.
  \emph{Negative branch}, by cases on the family's head
  (Lemma~\ref{lem:coreinv}; there is no $\Glue$ case and no neutral
  case):

  \emph{$\mathbb{N}$-line:} the family value at the generic
  dimension has head $\mathsf{v}\mathbb{N}$ with no arguments, so
  it is $\ell$-free and the check fired
  (Theorem~\ref{thm:exact}); this branch is unreachable.

  \emph{$\Pi$-line} $\langle i \rangle \Pi x{:}T.S$: the
  contractum is $\lambda x_1.\, \transp^i
  S[w\,x_1\,i/x]\,(f\,(w\,x_1\,0))$ with
  $w\,x_1\,i = \transp^j T(i \vee \neg j)\,x_1$.  Apply to related
  $x_1 \sim x_1'$ at the endpoint domain: $w\,x_1\,i$ is a
  transport along the \emph{domain line} --- the syntactic
  subexpression $T$ under a reparametrized interval environment
  --- so this lemma applies at $T$ (structural recursion,
  Remark~\ref{rem:order}), giving related fills, in particular
  related $w\,x_1\,0$ at the source domain; $f$ applied to them is
  related at $S[\rho{+}[w\,x_1\,0]]$ by the $\Pi$-clause; the
  outer transport is along $\langle i \rangle
  S[w\,x_1\,i/x]$ --- the subexpression $S$ in the environment
  extended by the fill --- so this lemma at $S$ closes the case,
  and the results assemble by the $\Pi$-clause at the endpoint
  type.

  \emph{$\Sigma$-line} $\langle i \rangle \Sigma x{:}T.S$: the
  first component transports along $T$ (this lemma at $T$).  The
  \emph{codomain typing detail}: the second component transports
  along $\langle i \rangle S[\overline{u}\,i/x]$ where
  $\overline{u}\,i = \transp^j T(i \wedge j)\,(\fst u)$ is the
  first component's fill; by this lemma at $T$ the fills are
  related at $T[\rho{+}[\iota]]$ for each $\iota$, so the
  environments $\rho{+}[\overline{u}\,i]$ are related, the line
  $\langle i \rangle S[\overline{u}\,i/x]$ is a computable family
  (its instantiations are evaluations of the subexpression $S$ in
  related environments), and this lemma at $S$ applies.  The pair
  is related at the endpoint type by the $\Sigma$-clause, whose
  second-component environment $\rho{+}[\transp\,T\,(\fst u)]$ is
  exactly the fill at $\ii$.

  \emph{$\mathsf{PathP}$-line}
  $\langle i \rangle \mathsf{PathP}(\langle j \rangle T)\,l\,r$
  (where $T, l, r$ may mention $i$): the contractum and its
  obligations, in derivation form.  The structural rule produces
  \[
    \langle j \rangle\, \vcomp^i\, T(i,j)\,
      [\,j{=}0 \mapsto l(i),\ j{=}1 \mapsto r(i)\,]\,(p\,j),
  \]
  whose premises are exactly the composition rule's:
  \begin{align}
    &j{=}0:\quad l(0) \conv p\,0
      && \text{the source typing of } p
      \tag{P1}\label{eq:P1}\\
    &j{=}1:\quad r(0) \conv p\,1
      && \text{likewise}
      \tag{P2}\label{eq:P2}\\
    &j{=}0 \wedge j{=}1:\quad \text{vacuous}
      && (0 = 1) = \bot
      \tag{P3}\label{eq:P3}
  \end{align}
  --- \eqref{eq:P1}/\eqref{eq:P2} hold because
  $p : \mathsf{PathP}(\langle j \rangle T(0,j))\,l(0)\,r(0)$ and
  path application at an endpoint returns that endpoint.  The
  composite unfolds as transport of the base and tubes along
  $\langle i \rangle T(i,j)$ followed by an $\hcomp$ at $T(1,j)$
  with the transported endpoint tubes; each transport is this lemma
  at the subterm $T$, and the $\hcomp$ is
  Lemma~\ref{lem:corehcompstr} at $T$.  The \emph{endpoint
  coherence} of the result --- required by the
  $\mathsf{PathP}$-clause at the target type
  $\mathsf{PathP}(\langle j \rangle T(1,j))\,l(1)\,r(1)$ --- is the
  face computation:
  \begin{align}
    (\text{result})\,\io
    &= \vcomp\text{-value on the decided face } j{=}0
    = \text{top of the } l\text{-tube} \notag\\
    &= \text{the } \transp^i\,T(i,0)\text{-corrected } l
       \text{ at } i{=}1
    \;\conv\; l(1),
    \tag{E0}\label{eq:E0}\\
    (\text{result})\,\ii
    &\conv r(1) \quad\text{symmetrically},
    \tag{E1}\label{eq:E1}
  \end{align}
  where the last step of \eqref{eq:E0} is the composition rule's
  face collapse ($\hcomp$ on a decided face returns the tube top)
  together with the tube's own transport landing at $l(1)$ ---
  precisely the boundary rule of $\vcomp$.  So the contractum
  inhabits the target path type with the required endpoints, and
  relatedness at generic $j$ is the two ingredients above.
\end{proof}

\subsection{The fundamental theorem, and the corollaries}
\label{app:corefund}

\begin{thm}\label{thm:corefund}
  Every rule of $T^{\mathrm{core}}_{\mathrm{alg}}$ preserves
  semantic typing (Definition~\ref{def:semtyping}, with
  $\valeq_{\evalop(T,\rho)}$ read as $\sim_{T[\rho,\rho']}$).
\end{thm}
\begin{proof}
  Induction on the typing derivation.  Variables: environment
  relatedness.  $\lambda$/application, pairing/projections, path
  abstraction/application: the matching clauses of
  Section~\ref{app:corerel}, with definedness carried by each
  clause.  $\mathbb{N}$-introductions: numerals.
  $\mathsf{natrec}$: by the $\mathbb{N}$-clause the scrutinee is a
  numeral (compute by an inner induction on the numeral, the step
  case using the $\Pi$-clause of the step premise) or a neutral
  (both sides form recursor neutrals with equal read-back ---
  motive by reify under its binder ---
  related by reflect at the motive instance).  $\transp$:
  Lemma~\ref{lem:coretransp}.  $\hcomp$:
  Lemmas~\ref{lem:corehcompN} and~\ref{lem:corehcompstr}.
  Conversion rule: Lemma~\ref{lem:typeconvsound} (the switch of
  type predicates), with Lemma~\ref{lem:coreconvcomp} and
  Lemma~\ref{lem:coreper} (environment/type irrelevance).
\end{proof}

\begin{cor}\label{cor:corepkg}
  For the minimal core, unconditionally: evaluation, read-back and
  conversion terminate on well-typed inputs; canonicity (closed
  naturals are numerals: the $\mathbb{N}$-clause, no closed
  neutrals, and no closed stuck composites by the closed-instance
  remark of Lemma~\ref{lem:corehcompN}); consistency; and the
  admissibility package of Theorem~\ref{thm:admissible}.
\end{cor}

\subsection{A paper proof of the path separation}
\label{app:papersep}

The read-back characterization of conversion
(Lemma~\ref{lem:convrb}) upgrades the machine witness of
Lemma~\ref{lem:pathfail} to a paper theorem, for every fixed core
type.

\begin{lem}[$\hcomp$/spine separation]\label{lem:sepschema}
  Let $T$ be a core type over a core context, $n$ a neutral of type
  $T$, and
  $H := \hcomp^{\,\iota}\,T\,[\vec\varphi \mapsto \vec t\,]\,(n)$ a
  composite whose faces $\vec\varphi$ are undecided at the current
  world and whose tubes are neutral-based values constant in the
  fill dimension $\iota$.  Then $\quo(H) \ne \quo(n)$; hence, by
  Lemma~\ref{lem:convrb}, $H \not\aconv n$.
\end{lem}
\begin{proof}
  Structural induction on $T$, following the kernel's structural
  $\hcomp$ rules (which are also how $\quo$ sees the values).
  \emph{$\mathbb{N}$:} the base is neutral, so $H$ is a stuck
  $\hcomp$-neutral (Section~\ref{app:coretypes}); its read-back is
  an $\hcomp$-term, while $\quo(n)$ is an elimination spine ---
  distinct productions of the normal-form grammar.
  \emph{$\Pi x{:}T_1.S$:} read-back $\eta$-expands both sides at
  the generic $x_d$; on the left, the structural rule turns
  $H\,x_d$ into an $\hcomp$ at $S[x_d]$ whose base $n\,x_d$ is
  again a neutral spine and whose tubes remain fill-constant and
  neutral-based; on the right, $n\,x_d$ is a spine.  The induction
  hypothesis at $S$ gives distinct body read-backs, hence distinct
  $\lambda$-bodies.
  \emph{$\Sigma x{:}T_1.S$:} read-back compares first projections;
  $\fst H$ is the composite of the first components over the spine
  $\fst n$ (fill-constant, neutral-based), and the induction
  hypothesis at $T_1$ separates already the first components.
  \emph{$\mathsf{PathP}(\langle k \rangle T_1)\,l\,r$:} read-back
  instantiates at the generic level $k_d$; the structural rule
  extends the system with the endpoint tubes $k{=}0 \mapsto l$,
  $k{=}1 \mapsto r$ --- undecided at generic $k_d$, so the
  composite does not collapse --- with base $n\,k_d$, a spine; the
  induction hypothesis at $T_1$ applies with the enlarged (still
  fill-constant, neutral-based) system.
  In every case the faces stay undecided at the generic
  instantiations, so no collapse intervenes, and the separation
  bottoms out at $\mathbb{N}$ in distinct grammar productions.
\end{proof}

\begin{prop}[Path separation in the core, on paper]
  \label{prop:papersep}
  For every core type $A$ and the generic context
  $(a\,b : A,\ p : \Path\,A\,a\,b)$:
  \[
    \langle j \rangle\, \hcomp^{\,\iota} A\,
      [\,j{=}0 \mapsto a,\ j{=}1 \mapsto b\,]\,(p\,j)
    \;\not\aconv\; p ,
  \]
  unconditionally.
\end{prop}
\begin{proof}
  Conversion at the path type instantiates both sides at the
  generic level $j_d$: the faces $j_d{=}0$, $j_d{=}1$ are
  undecided, the base $p\,j_d$ is a neutral spine, and the tubes
  $a, b$ (values of the context variables) are generic neutrals,
  constant in the fill dimension.  Lemma~\ref{lem:sepschema} at $A$
  separates the two sides; Lemma~\ref{lem:convrb} turns the
  read-back difference into rejection.  Every step lives in the
  minimal core, so Corollary~\ref{cor:corepkg} makes the statement
  unconditional.  Consequently the no-go analysis of
  Section~\ref{sec:nogo} --- both the derivation of
  \eqref{eq:dagger} in $T_{\mathrm{eq}}$ and its refutation by the
  algorithm --- is established on paper within the unconditional
  layer; the machine probes witness, in addition, the full
  implemented calculus with $A$ a universe variable.
\end{proof}

\subsection{The derived $\J$, explicitly}
\label{app:coreJ}

\begin{prop}[Strict $\J$ in the minimal core; a schema]\label{prop:J}
  Work in the core over a context
  $\Gamma = (x : A;\ \text{and, in the extended context } y : A,\
  p : \Path\,A\,x\,y,\ \text{a core type } C)$, with
  $d : C[x, \refl_x]$.  Define, as usual in cubical type theory,
  \[
    \J_C\,d\,y\,p \;:=\;
    \transp^i\, C\bigl[\,p(i),\ \langle j \rangle\, p(i \wedge j)
    \,\bigr]\ \io\ d .
  \]
  Then: (i) $\J_C$ is well-typed in the core (the family is a core
  type in the interval binder $i$, using the connection
  $i \wedge j$; no $\Glue$, universe or HIT is involved);
  (ii) $\J_C\,d\,x\,\refl_x \aconv d$, unconditionally.
\end{prop}
\begin{proof}
  (i) is the standard typing of the connection square: at $i=0$
  the family is $C[p(0), \langle j \rangle p(0)] =
  C[x, \refl_x]$ by the boundary and connection laws ($0 \wedge j
  = 0$), so $d$ is accepted at the source.
  (ii) Substitute $p := \refl_x$ and evaluate the transport's
  family at the fresh dimension $\ell$:
  $p(\ell)$ evaluates to (the value of) $x$, and
  $\langle j \rangle\, p(\ell \wedge j)$ evaluates to the constant
  path abstraction over $x$ --- the connection is absorbed by the
  DNF normalization ($\ell \wedge j$ instantiated on the constant
  path collapses; machine witness \texttt{strictConnZeroD}).  So
  the family value is the value of $C[x, \refl_x]$ in an
  environment not involving $\ell$, and by Lemma~L1 its read-back
  does not contain $\ell$.  By Theorem~\ref{thm:exact} the check
  fires; the positive branch returns the value of $d$; by
  Lemma~\ref{lem:coreconvcomp} the algorithmic equality
  $\J_C\,d\,x\,\refl_x \aconv d$ holds.  Every ingredient ---
  evaluation, check, conversion --- is inside the minimal core, so
  Corollary~\ref{cor:corepkg} applies and no assumption is used.
  Machine witness: \texttt{strictJReflD}.
\end{proof}

\begin{rem}[Schema, not internal term]
  The core has no universe, so $\J_C$ is a \emph{meta-level schema}:
  one derived eliminator per core type $C$ in the extended context,
  not a single internal universe-polymorphic term.  This is the
  usual status of $\J$ in a universe-free fragment and involves no
  loss: Proposition~\ref{prop:J} is proved uniformly in $C$.  (In
  the full calculus, where universes exist, the internal $\J$ is the
  same construction at a universe-valued motive; it then lives in
  the relative layer.)
\end{rem}

\subsection{The $S^1$ extension}
\label{app:s1ext}

\begin{prop}[$S^1$, proof sketch]\label{prop:s1}
  The results of this appendix extend to the fragment enlarged by
  $S^1$ ($\mathsf{base}$, $\mathsf{loop}$, its eliminator, and
  $\hcomp$-cells).
\end{prop}
\begin{proof}[Proof sketch (deliberately so labelled)]
  The $S^1$-clause of the relation adds: both $\mathsf{base}$;
  or $\mathsf{loop}$ at DNF-equal interval values; or
  componentwise-related $\hcomp$-cells with equal cofibrations; or
  neutrals with equal read-back.  Transport along a constant
  $S^1$-line always takes the positive branch (as for
  $\mathbb{N}$).  The eliminator computes on $\mathsf{base}$ and
  $\mathsf{loop}$ by its two rules, acts on an $\hcomp$-cell by
  composing in the motive fiber against the tube-transformed
  system (the kernel's tube-transformer closure), and is neutral
  on neutrals; boundary computation
  ($\mathsf{loop}(\io) = \mathsf{loop}(\ii) = \mathsf{base}$) is
  the face collapse.  Filling in the eliminator-over-$\hcomp$ case
  at the level of detail of
  Sections~\ref{app:corehcomp}--\ref{app:coretransp} is routine
  but lengthy, and we deliberately keep it \emph{outside} the
  unconditional theorem; the artifact provides machine evidence
  (the $S^1$ probes, the winding-number development
  $\pi_1(S^1) \cong \mathbb{Z}$, and the switchover probes for
  higher-inductive constructors).
\end{proof}

\subparagraph*{What the core deliberately omits.}  $\Glue$ (hence
$\ua$ and $\mathsf{transpGlue}$), universes (hence large
elimination), and all HITs (with $S^1$ recovered as
Proposition~\ref{prop:s1}).  These are exactly the ingredients
whose base metatheory is the content of
\cite{HuberNorm,SterlingAngiuli}; extending the discharge to them
is the remaining gap between Theorem A$'$ and Theorems A/B, and we
prefer to leave it visibly open rather than to compress those
proofs beyond inspectability.

\section{The Sterling--Angiuli bridge}
\label{app:bridge}

This appendix makes precise which parts of the base component (B) are
discharged by the normalization theorem of Sterling--Angiuli
\cite{SterlingAngiuli} for the Cartesian layer (Theorem A), and which
parts remain a transposition obligation.  It is based on the
published LICS 2021 paper (arXiv:2101.11479v2).

\subsection{What their theorem provides}

Their subject is intensional \emph{Cartesian} cubical type theory,
presented as the locally Cartesian closed category of judgments
generated by an explicit signature: connectives $\Pi$, $\Sigma$,
$\mathsf{path}$, $\mathsf{glue}$ (with $\mathsf{Equiv}$ defined
without universes), the circle $\mathsf{S1}$ with its eliminator, and
Kan operations $\mathsf{hcom}$ and
\[
  \mathsf{coe} : \textstyle\prod_{A : \I \to \mathsf{tp}}
  \prod_{r, s : \I} \mathsf{tm}(A(r)) \to
  \{\mathsf{tm}(A(s)) \mid r = s \hookrightarrow a\},
\]
with cofibration classifier $\mathbb{F}$ closed under
$\wedge, \vee, \forall_\I, =_\I$.  The method is synthetic Tait
computability --- a gluing model over the generic model, with the
novel device of \emph{stabilized neutrals} (neutrals indexed by the
cofibration recording where they cease to be neutral); it is, by
design, \emph{not} evaluation-based.  The main results are: the
normalization theorem (their Theorem~42), a bijection between
equivalence classes of terms in context and $\beta/\eta$-normal
forms; injectivity of type constructors (their Theorem~43); and an
algorithm deciding judgmental equality (their Corollary~47).  Two
scope restrictions are explicit in their paper: \emph{universes are
not treated} (they note the expected extension via cumulative
universes and the tininess of the interval, but do not carry it
out), and the only higher inductive type is the circle.  For the
De~Morgan variant they state that the argument ``carries over
mutatis mutandis''; the adaptation is likewise not carried out.

\subsection{Item-by-item correspondence}

\begin{center}
\begin{tabular}{p{0.44\textwidth}p{0.50\textwidth}}
  component (B) item & status under \cite{SterlingAngiuli} \\\hline
  normal-form language; uniqueness; inconvertibility of distinct
  normal forms (L13, second half) &
  \textbf{provided}: Theorem 42 (bijection with $\beta/\eta$-normal
  forms), for $\Pi, \Sigma, \mathsf{path}, \mathsf{glue},
  \mathsf{S1}$, no universes \\
  injectivity of type constructors &
  \textbf{provided}: Theorem 43 \\
  decidability of the base judgmental equality &
  \textbf{provided}: Corollary 47 \\
  predicate clauses of $\valeq$ for their connectives &
  \textbf{provided in substance}: their computability structures
  (with stabilized neutrals) are the proof-relevant counterpart of
  Definition~\ref{def:valeq}'s clauses; the transposition to a
  concrete Kripke PER over our domain is mechanical but must be
  written \\
  environment-form fundamental cases for a concrete NbE evaluator
  (Appendix~\ref{app:fund}, (B)-cases); read-back correctness for
  that evaluator (L13, first half); evaluator termination content &
  \textbf{not literal}: their normalization function arises from the
  glued model, not from an operational evaluator; the standard
  NbE-correctness argument for a concrete evaluator against their
  normal forms is a known pattern but is not in their paper \\
  universes ($\valeq_\U$, universe rules of
  Appendix~\ref{app:fund}) & \textbf{gap}: outside their theorem;
  expected but unproven extension \\
  HITs beyond the circle (Section~\ref{sec:hits}) &
  \textbf{gap}: outside their theorem \\
  De~Morgan instance & \textbf{gap}: ``mutatis mutandis'' claim, not
  carried out --- which is why our De~Morgan layer is Theorem B,
  relative to (B) \\
\end{tabular}
\end{center}

\subsection{Primitive translation: boundary-conditioned
  $\mathsf{coe}$ vs.\ $\varphi$-annotated $\transp$}

Their coercion is regular \emph{on the diagonal}: the codomain
$\{\mathsf{tm}(A(s)) \mid r = s \hookrightarrow a\}$ builds
$\mathsf{coe}(A, r, r, a) = a$ into the \emph{type} of the primitive
--- the Cartesian counterpart of our $(\mathrm{S}_{\ii})$, expressed
as a boundary condition rather than a formula annotation.  Our
$\transp^i A\,\varphi\,u$ translates as
$\mathsf{coe}(\lambda i.A, 0, 1, u)$, with the constancy formula
$\varphi$ corresponding to the cofibration under which the line is
required constant; rule (R) then reads: coerce along a line whose
\emph{normal form} is degenerate as if $r = s$.  The switchover
analysis of Section~\ref{sec:switchover} is insensitive to the
difference: the derivation \eqref{eq:d1}--\eqref{eq:dagger} uses only
the diagonal rule, the structural path rule and congruence, all of
which the Cartesian signature has, and the residue is an
$\mathsf{hcom}$ with fill-constant tubes.

\subsection{The precise statement of Theorem A}

\begin{quote}
  \textbf{Theorem A.}  Over the fragment
  $\Pi, \Sigma, \mathsf{path}, \mathsf{glue}, \mathsf{S1}$ of
  Cartesian cubical type theory --- the fragment covered by
  \cite{SterlingAngiuli} --- the results of
  Sections~\ref{sec:oper}--\ref{sec:canonicity}
  (admissibility, termination, canonicity, definitional
  $\J\,d\,\refl$) hold for the prioritized operational system,
  \emph{modulo the evaluator-level transposition}: the re-proof, for
  a concrete NbE evaluator of that fragment, of the (B)-marked
  lemmas of Appendices~\ref{app:cases}--\ref{app:fund} against the
  normal forms of \cite{SterlingAngiuli}.  The transposition is
  standard (it is the routine half of every NbE-correctness proof)
  and involves no cubical subtlety beyond what
  \cite{SterlingAngiuli} already resolved, but it is a proof
  obligation, not a citation; and universes, the wider HIT roster,
  and the De~Morgan instance remain outside it.
\end{quote}

\noindent We prefer this honest statement to the stronger
``unconditional'' phrasing: the alternative --- claiming the
evaluator-level lemmas from a theorem that is deliberately not
evaluation-based --- would misattribute the one part of the base
component that their method, precisely by its elegance, does not
supply.

\begin{thebibliography}{10}

\bibitem{CCHM}
C.~Cohen, T.~Coquand, S.~Huber, A.~M\"ortberg.
\newblock Cubical type theory: a constructive interpretation of the
  univalence axiom.
\newblock \emph{TYPES 2015}, LIPIcs 69, 2018.

\bibitem{ABCFHL}
C.~Angiuli, G.~Brunerie, T.~Coquand, R.~Harper, K.-B.~Hou, D.~Licata.
\newblock Syntax and models of Cartesian cubical type theory.
\newblock \emph{Math.\ Struct.\ Comput.\ Sci.}, 31(4), 2021.

\bibitem{VezzosiMortbergAbel}
A.~Vezzosi, A.~M\"ortberg, A.~Abel.
\newblock Cubical Agda: a dependently typed programming language with
  univalence and higher inductive types.
\newblock \emph{J.\ Funct.\ Program.}, 31, 2021.

\bibitem{HuberNorm}
S.~Huber.
\newblock Canonicity for cubical type theory.
\newblock \emph{J.\ Autom.\ Reasoning}, 63, 2019.

\bibitem{SterlingAngiuli}
J.~Sterling, C.~Angiuli.
\newblock Normalization for cubical type theory.
\newblock \emph{LICS 2021}.

\bibitem{SwanRegularity}
A.~W.~Swan.
\newblock Separating path and identity types in presheaf models of
  univalent type theory.
\newblock Preprint, arXiv:1808.00920, 2018.

\bibitem{SattlerNote}
C.~Sattler.
\newblock The failure of regularity for composition in the CCHM
  universe (unpublished observation, commonly attributed; we know of
  no self-contained publication).
\newblock Published accounts: the introduction and historical
  remarks of \cite{SwanRegularity}, and the discussion of regular
  composition in \cite{CCHM,ABCFHL}.

\bibitem{Arend}
V.~Isaev, F.~Part, S.~Sinchuk.
\newblock Arend: a theorem prover based on homotopy type theory.
\newblock JetBrains Research.  Language paper and the Prelude
  documentation of \textsf{coe} (whose computation rules include the
  constant-family collapse) at
  \url{https://arend-lang.github.io/documentation/language-reference/prelude};
  accessed 2026-07-17.

\bibitem{UIPCubical}
Y.~J.~Tan.
\newblock Towards computational UIP in Cubical Agda.
\newblock Preprint, arXiv:2511.21209, 2025.

\bibitem{HuangNF}
X.~Huang.
\newblock Normal forms in cubical type theory.
\newblock Preprint, arXiv:2603.24923, 2026.

\bibitem{OrtonPitts}
I.~Orton, A.~M.~Pitts.
\newblock Axioms for modelling cubical type theory in a topos.
\newblock \emph{LMCS} 14(4), 2018.

\bibitem{StickyXTT}
J.~Sterling, C.~Angiuli, D.~Gratzer.
\newblock Cubical syntax for reflection-free extensional equality.
\newblock \emph{FSCD 2019}.

\end{thebibliography}

\end{document}
